%% ============================================================================
%% Chapter 18: Observability & SRE for APIs
%% Mastering APIs: From Zero to Production Guru
%% ============================================================================
\chapter{Observability \& SRE for APIs}
\label{ch:observability}

%% ----------------------------------------------------------------------------
%% Executive Summary
%% ----------------------------------------------------------------------------
Every chapter until now has asked a design-time question: how should the
API be shaped, secured, versioned, and deployed? This chapter asks the
runtime question that decides whether all of that work matters: \emph{when
the API misbehaves in production --- and it will --- can you see it,
measure it, and respond before your error budget is gone?} An API without
observability is a system whose failures are reported by customers, whose
latency regressions are discovered in quarterly reviews, and whose
outages are debugged by intuition. An API without SLOs has observability
data but no judgment: dashboards nobody pages on, alerts everybody ignores,
and no principled answer to ``is this bad enough to stop the release?''

We build the answer in three layers. First, the \textbf{three pillars}
applied concretely to APIs: structured JSON logs with levels, sampling,
and a hard rule against PII and secrets; metrics organized as RED (rate,
errors, duration) for services and USE (utilization, saturation, errors)
for resources; and distributed traces with spans and context propagation.
Second, the \textbf{machinery}: OpenTelemetry's API/SDK/exporter/collector
architecture~\cite{opentelemetrydocs}, the W3C \texttt{traceparent} header
parsed byte by byte (version, trace-id, span-id, flags), \texttt{tracestate}
and baggage, span kinds and HTTP semantic conventions, and the head-versus-tail
sampling decision. Third, the \textbf{engineering of numbers}: histogram
buckets and why percentiles must come from histograms (including a proof
by example of the p99-of-averages fallacy), cardinality discipline with
cost arithmetic, exemplars, SLI formulas, error budgets (43.2 minutes per
month at 99.9\%), and multi-window multi-burn-rate alerting with the full
window-pair mathematics~\cite{beyer2016sre}.

The operational half closes the loop: correlation IDs end to end
(connecting the problem-details work of Chapter~\ref{ch:problem_details}),
per-endpoint dashboards anchored to SLOs rather than vibes, a systematic
timeline-decomposition method for debugging production latency, and the
incident-response apparatus --- severity levels, roles, runbooks, and
blameless postmortems with a template. All code is typed Python~3.11,
offline-first (the OTLP exporter is replaced by a file exporter; the
collector is simulated), deterministic (seeded ID generators and
timelines), and set in the Northwind Robotics fleet-telemetry universe.
The war story is the invisible retry storm: a latency regression that
multiplied traffic fivefold while every individual log line looked
perfectly healthy --- found through trace fan-out metrics, invisible in
logs.

\begin{keyconceptbox}
Observability answers \emph{``what is the system doing?''}; SRE answers
\emph{``is that good enough, and what do we do when it is not?''} The two
are useless apart: telemetry without SLOs is data hoarding, and SLOs
without telemetry are wishes. The contract that binds them is the error
budget --- the only number in this chapter that is allowed to page a human
being --- and the mechanism that makes telemetry coherent across service
boundaries is context propagation via the W3C \texttt{traceparent} header.
\end{keyconceptbox}

%% ----------------------------------------------------------------------------
\section{Learning Objectives \& Prerequisites}
\label{sec:ch18-objectives}

\subsection*{Learning Objectives}
After completing this chapter, its lab, and its exercises, you will be able
to:

\begin{enumerate}[label=\textbf{LO-\arabic*.}, leftmargin=2.6cm]
  \item Apply the three pillars to API services concretely: emit structured
        JSON logs with levels and sampling and no PII or secrets; define RED
        metrics for services and USE metrics for resources; construct traces
        with correct span kinds and context propagation.
  \item Explain the OpenTelemetry architecture --- API, SDK, exporter,
        collector --- and justify the collector as the boundary between
        application code and vendor backends~\cite{opentelemetrydocs}.
  \item Parse and emit the W3C \texttt{traceparent} header by hand
        (version--trace-id--span-id--flags), explain \texttt{tracestate}
        and baggage, and state exactly why context does \emph{not} propagate
        itself through queues.
  \item Engineer latency metrics correctly: choose histogram buckets,
        compute percentiles from histograms, refute the p99-of-averages
        fallacy, budget label cardinality with explicit arithmetic, and use
        exemplars to jump from a metric to a trace.
  \item Design SLIs for availability, latency, and freshness; derive error
        budgets (43.2 minutes per month at 99.9\%); and implement
        multi-window multi-burn-rate alerting with the fast/slow window-pair
        mathematics of the SRE Workbook~\cite{beyer2016sre}.
  \item Run operational practice end to end: correlation IDs propagated to
        problem-details responses, SLO-anchored per-endpoint dashboards,
        timeline decomposition of a slow request, and incident response with
        severities, roles, runbooks, and blameless postmortems.
  \item Instrument a FastAPI service with traces, metrics, and correlated
        logs offline, and verify propagation with executable
        tests~\cite{fastapidocs}.
\end{enumerate}

\subsection*{Prerequisites}
This chapter assumes Chapters~0--17, especially: building and testing
FastAPI services (Chapter~\ref{ch:fastapi}), structured errors and
correlation-friendly problem details (Chapter~\ref{ch:problem_details}),
retries and resilience patterns (Chapter~\ref{ch:microservices}), API
security boundaries --- you cannot judge what may be logged without them
(Chapter~\ref{ch:api_security}), and deployment mechanics
(Chapter~\ref{ch:deployment}). Elementary probability (quantiles, ratios)
and algebra are assumed; every formula is derived in place.

\subsection*{Setup}
All listings run offline on Windows with Python~3.11+. One-time setup from
PowerShell in a scratch directory:

\begin{minted}{text}
python -m venv .venv
.\.venv\Scripts\Activate.ps1
pip install "opentelemetry-sdk>=1.25" `
            "opentelemetry-instrumentation-fastapi>=0.46b0" `
            "fastapi>=0.110" "pytest>=8"
\end{minted}

No collector, no network, no credentials: spans are exported to a local
JSON-lines file, metrics are read back from an in-memory reader, and every
timeline is seeded. Section~\ref{sec:ch18-deps} notes what changes when you
point the same code at a real OpenTelemetry Collector.

%% ----------------------------------------------------------------------------
\section{Deep Technical Mechanics \& Architectural Concepts}
\label{sec:ch18-mechanics}

\subsection{The Three Pillars, Applied to APIs}
\label{subsec:ch18-pillars}

``Logs, metrics, traces'' is a slogan; the engineering content is in what
each pillar is \emph{for}, what it \emph{costs}, and how they \emph{join}.

\paragraph{Structured logs: events with a schema.} A log line is a
timestamped record of something that happened once. The moment an API
serves real traffic, free-text logs stop being searchable and start being
forensics-resistant, so the discipline is: emit JSON, one object per line,
with a fixed schema --- timestamp, level, logger, message, and
\emph{correlation fields} (trace id, span id). Levels are a contract with
the reader: \texttt{ERROR} means a human should care, \texttt{WARN} means a
human should eventually care, \texttt{INFO} means normal operation worth
keeping, \texttt{DEBUG} means sampled or off in production. Volume control
is part of the design: a hot endpoint logging every request at
\texttt{INFO} is a denial-of-wallet attack on yourself, so high-frequency
paths are sampled (keep 1 in $N$) while errors are always kept. And the
hard rule, discussed again in Section~\ref{sec:ch18-pitfalls}: \textbf{no
PII, no credentials, no tokens, no request bodies} --- logs outlive their
retention assumptions, leak into support tickets and vendor UIs, and are
readable by people who should not see user data.

\paragraph{Metrics: aggregatable numbers.} A metric is a number observed
repeatedly and aggregated cheaply: counters (requests served), histograms
(durations in buckets), gauges (queue depth). For API services the working
vocabulary is \textbf{RED}: \emph{Rate} (requests per second),
\emph{Errors} (failures per second, or error ratio), \emph{Duration}
(latency distribution). For the resources under the service, the
complementary vocabulary is \textbf{USE}: \emph{Utilization} (how busy),
\emph{Saturation} (how queued), \emph{Errors} (how broken) --- connection
pools, worker processes, disk. RED tells you your users are hurting; USE
usually tells you why.

\paragraph{Traces: causality with timestamps.} A trace records one request's
journey as a tree of \emph{spans} --- named, timed intervals with a parent
pointer. Where logs answer ``what happened here?'' and metrics answer ``how
often, how fast?'', traces answer the question microservices make urgent:
``\emph{where} did the time go, across the three services this request
touched?'' The trace/span model and its wire format get their own
subsections below.

\begin{definition}[Span, Trace, and Context Propagation]
A \emph{span} is a named, timed operation with attributes, a status, a
span id, and an optional parent span id. A \emph{trace} is the set of spans
sharing one trace id, forming a tree. \emph{Context propagation} is the
mechanism by which a caller transmits the trace id, its own span id, and
sampling flags to a callee --- for HTTP, via the W3C \texttt{traceparent}
header --- so the callee's spans attach to the same tree.
\end{definition}

\begin{keyconceptbox}
The pillars join on \textbf{exemplars and correlation ids}: a metric
histogram bucket carries an exemplar pointing at one trace; every log line
emitted inside a span carries that span's trace id; every problem-details
response (Chapter~\ref{ch:problem_details}) carries the same id for the
client to quote back. One identifier, three pillars, zero archaeology.
\end{keyconceptbox}

\subsection{OpenTelemetry Architecture: API, SDK, Exporter, Collector}
\label{subsec:ch18-otel}

OpenTelemetry (OTel) is the industry-converged instrumentation standard,
and its layering is the first architecture decision to get
right~\cite{opentelemetrydocs}:

\begin{itemize}
  \item \textbf{API} --- the instrumentation surface your code calls:
        \texttt{tracer.start\_as\_current\_span(...)},
        \texttt{meter.create\_counter(...)}. Depending on the API is free;
        it does nothing until an SDK is installed.
  \item \textbf{SDK} --- the implementation: span processors, samplers,
        metric aggregations, resource attribution (\texttt{service.name}).
        Applications configure the SDK; libraries must never.
  \item \textbf{Exporter} --- the egress adapter: OTLP over gRPC/HTTP to a
        collector in production; in this chapter, a deliberately offline
        JSON-lines file exporter implementing the same interface, so the
        telemetry pipeline is testable with zero infrastructure.
  \item \textbf{Collector} --- a standalone process receiving OTLP,
        processing (batching, redaction, tail sampling), and re-exporting
        to backends (Prometheus, Jaeger/Tempo, Loki/Elasticsearch, or a
        vendor). Putting the collector between app and backend is what lets
        you switch vendors by editing one YAML file instead of a fleet
        redeploy.
\end{itemize}

\begin{figure}[ht]
\centering
\begin{tikzpicture}[
  font=\small, >=stealth,
  flowstep/.append style={minimum height=9mm}
]
  % ---- app tier ----
  \node[flowstep, fill=blue!8] (appcode) {Application code\\[-2pt]\footnotesize OTel API calls};
  \node[flowstep, fill=blue!8, right=6mm of appcode] (sdk) {OTel SDK\\[-2pt]\footnotesize processors, sampler};
  \node[flowstep, fill=blue!8, right=6mm of sdk] (exporter) {Exporter\\[-2pt]\footnotesize OTLP (or file)};
  \node[draw, rounded corners, dashed, inner sep=5pt, fit=(appcode)(sdk)(exporter),
        label={[modcap]above:Application process}] (app) {};
  % ---- collector ----
  \node[flowstep, fill=packtorange!12, right=14mm of exporter] (recv) {Receivers\\[-2pt]\footnotesize OTLP gRPC/HTTP};
  \node[flowstep, fill=packtorange!12, right=6mm of recv] (proc) {Processors\\[-2pt]\footnotesize batch, redact, tail-sample};
  \node[flowstep, fill=packtorange!12, right=6mm of proc] (colexp) {Exporters\\[-2pt]\footnotesize per-backend};
  \node[draw, rounded corners, dashed, inner sep=5pt, fit=(recv)(proc)(colexp),
        label={[modcap]above:OpenTelemetry Collector}] (coll) {};
  % ---- backends ----
  \node[flowstep, fill=green!10, right=14mm of colexp, yshift=10mm] (prom) {Prometheus\\[-2pt]\footnotesize metrics};
  \node[flowstep, fill=green!10, right=14mm of colexp] (jaeger) {Jaeger / Tempo\\[-2pt]\footnotesize traces};
  \node[flowstep, fill=green!10, right=14mm of colexp, yshift=-10mm] (loki) {Loki / ES\\[-2pt]\footnotesize logs};
  \node[draw, rounded corners, dashed, inner sep=5pt, fit=(prom)(loki),
        label={[modcap]right:Backends}] {};
  % ---- edges ----
  \draw[->, thick] (appcode) -- (sdk);
  \draw[->, thick] (sdk) -- (exporter);
  \draw[->, thick] (exporter) -- (recv);
  \draw[->, thick] (recv) -- (proc);
  \draw[->, thick] (proc) -- (colexp);
  \draw[->, thick] (colexp.east) -- (prom.west);
  \draw[->, thick] (colexp.east) -- (jaeger.west);
  \draw[->, thick] (colexp.east) -- (loki.west);
  \node[note, below=2mm of app] {library + app instrumentation};
  \node[note, below=2mm of coll] {vendor-neutral boundary};
\end{tikzpicture}
\caption{The observability pipeline. Application code depends only on the
OTel API; the SDK assembles spans and metric aggregations; the exporter
ships OTLP to the collector, which batches, redacts, optionally
tail-samples, and fans out to one backend per signal. Vendor changes happen
at the collector, never in application code.}
\label{fig:ch18-pipeline}
\end{figure}

Two SDK details matter for API services. First, \textbf{span processors}:
the production default is the \emph{batch} processor (queue and flush every
few seconds; the hot path only enqueues); this chapter's offline listings
use the \emph{simple} processor (export synchronously on span end) because
deterministic tests beat throughput here. Second, \textbf{metric views}:
the SDK lets you re-shape an instrument's aggregation --- we use a view to
pin the latency histogram to explicit millisecond buckets, which is what
makes percentile math and backend cost predictable
(Section~\ref{subsec:ch18-metrics-engineering}).

\subsection{The W3C Trace Context: \texttt{traceparent}, Byte by Byte}
\label{subsec:ch18-traceparent}

Context propagation across HTTP services is standardized by the W3C Trace
Context recommendation. The \texttt{traceparent} header is exactly four
dash-separated lowercase-hex fields:

\begin{examplebox}
\textbf{Decoding a live header.}
\texttt{00-4bf92f3577b34da6a3ce929d0e0e4736-00f067aa0ba902b7-01} means:
version \texttt{00}; trace id \texttt{4bf92f3577b34da6a3ce929d0e0e4736}
(128 bits); parent span id \texttt{00f067aa0ba902b7} (64 bits); flags
\texttt{01} --- bit~0 set, so this trace is \emph{sampled}. Listing~18.1
parses this format strictly and Listing~18.2 proves the round trip.
\end{examplebox}

\begin{table}[ht]
  \centering
  \small
  \begin{tabularx}{\textwidth}{l l l X}
    \toprule
    \textbf{Field} & \textbf{Width} & \textbf{Example} & \textbf{Rules} \\
    \midrule
    \texttt{version}   & 2 hex  & \texttt{00} & \texttt{ff} is forbidden; version \texttt{00} forbids trailing fields; future versions may append fields, which parsers must tolerate. \\
    \texttt{trace-id}  & 32 hex & \texttt{4bf9\ldots4736} & All-zero is invalid. Identifies the whole trace; constant across every hop. \\
    \texttt{span-id}   & 16 hex & \texttt{00f067aa0ba902b7} & All-zero is invalid. Identifies the caller's span; the callee records it as its \emph{parent}. \\
    \texttt{flags}     & 2 hex  & \texttt{01} & Only bit 0 (\texttt{0x01}, \emph{sampled}) is defined; other bits are reserved and must not be echoed as set. \\
    \bottomrule
  \end{tabularx}
  \caption{\texttt{traceparent} field reference (W3C Trace Context).}
  \label{tab:ch18-traceparent}
\end{table}

Three companion mechanisms complete the picture. \textbf{\texttt{tracestate}}
is a comma-separated list of at most 32 vendor \texttt{key=value} members
carrying vendor-specific sampling or routing state; it is mutable hop by
hop under strict size and syntax rules. \textbf{Baggage} is a separate W3C
header carrying user-defined key/value pairs (e.g.\ tenant or plan tier)
that propagate with the trace --- powerful, dangerous (it travels
\emph{everywhere}, so never put PII in it), and intentionally not part of
the span identity. And \textbf{queues}: nothing about an HTTP header helps
you across a message broker. Propagation through a queue requires
\emph{explicitly} writing the trace context into the message metadata on
publish and extracting it on consume (the same inject/extract pair,
different carrier) --- the async patterns of Chapter~\ref{ch:async_apis}
do not do this for you. Section~\ref{sec:ch18-pitfalls} lists the broken
trace tree as a first-class anti-pattern.

\subsection{Span Kinds and Semantic Conventions for HTTP}
\label{subsec:ch18-semconv}

OTel assigns every span a \emph{kind} that states its relationship to its
parent, and getting kinds right is what makes waterfall views and
service-graph aggregations correct~\cite{opentelemetrydocs}:

\begin{itemize}
  \item \textbf{SERVER} --- entry point of a remote call handled by this
        process (an inbound HTTP request). Parent came from headers.
  \item \textbf{CLIENT} --- an outgoing remote call (an HTTP request this
        process makes). Its span id becomes the child's parent via the
        injected \texttt{traceparent}.
  \item \textbf{INTERNAL} --- work inside the process (a cache lookup, a
        pricing calculation).
  \item \textbf{PRODUCER} / \textbf{CONSUMER} --- asynchronous messaging;
        the consumer's parent may not be on the critical path of the
        producer's latency, which is why they are distinct kinds.
\end{itemize}

Semantic conventions then standardize attribute names so dashboards and
backends are portable: \texttt{http.request.method},
\texttt{url.path} (templated route, never the raw path with ids),
\texttt{http.response.status\_code}, \texttt{server.address},
\texttt{http.route}. The single most consequential rule for APIs:
\textbf{attributes use route templates} (\texttt{/fleet/\{id\}}), never
concrete paths --- a concrete path per request is a cardinality bomb, the
subject of Section~\ref{subsec:ch18-metrics-engineering}.

\begin{figure}[ht]
\centering
\begin{tikzpicture}[font=\small, >=stealth]
  % ---- service swimlanes ----
  \node[modcap] at (0,0) {gateway};
  \node[modcap] at (5.2,0) {orders};
  \node[modcap] at (10.4,0) {inventory};
  \draw[gray!50, dashed] (0,-0.4) -- (0,-4.1);
  \draw[gray!50, dashed] (5.2,-0.4) -- (5.2,-4.1);
  \draw[gray!50, dashed] (10.4,-0.4) -- (10.4,-4.1);
  \node[note] at (0,-4.35) {time $\rightarrow$};
  % ---- spans as horizontal bars ----
  % gateway server span: t=0..48
  \fill[darkblue!75] (-0.06,-0.8) rectangle (4.86,-1.05);
  \node[note, anchor=west, text=darkblue] at (5.0,-0.92)
    {\texttt{gateway POST /fleet/telemetry} (SERVER, 48\,ms)};
  % orders server span: t=2..40
  \fill[packtorange!80] (5.14,-1.7) rectangle (8.96,-1.95);
  \node[note, anchor=west, text=packtorange] at (9.1,-1.82)
    {\texttt{orders.create} (SERVER, 38\,ms)};
  % inventory server span: t=8..30
  \fill[green!55!black] (10.34,-2.6) rectangle (12.54,-2.85);
  \node[note, anchor=west, text=green!50!black] at (8.6,-3.35)
    {\texttt{inventory.reserve} (SERVER, 22\,ms)};
  % ---- parent arrows ----
  \draw[->, thick, gray] (4.9,-0.92) .. controls (5.05,-1.2) .. (5.2,-1.7)
    node[midway, left, note] {parent};
  \draw[->, thick, gray] (9.0,-1.82) .. controls (9.9,-2.1) .. (10.3,-2.6)
    node[midway, below right=-2pt, note] {parent};
  % ---- traceparent hops ----
  \node[note, align=left] at (2.4,-3.9)
    {hop 1 inject: \texttt{traceparent: 00-4bf9\ldots4736-\emph{s1}-01}};
  \node[note, align=left] at (7.6,-3.9)
    {hop 2 inject: \texttt{00-4bf9\ldots4736-\emph{s2}-01}};
  \draw[->, thick, darkblue!60] (2.2,-3.55) -- (4.6,-2.1);
  \draw[->, thick, darkblue!60] (9.4,-3.55) -- (10.2,-3.0);
\end{tikzpicture}
\caption{One trace across three services. Bars are spans on a shared time
axis; every span carries trace id \texttt{4bf9\ldots4736}; each hop injects
a fresh span id into the outgoing \texttt{traceparent}, which the next
service records as its parent. This is exactly the tree Listing~18.3
exports and asserts.}
\label{fig:ch18-trace-tree}
\end{figure}

\subsection{Sampling: Head, Tail, and the Errors You Cannot Afford to Lose}
\label{subsec:ch18-sampling}

At 10{,}000 requests per second, keeping every trace is a storage and
egress decision most platforms lose. Sampling reduces volume; where the
decision happens determines what you keep.

\textbf{Head sampling} decides at the trace's first span: flip a weighted
coin (keep $p$\%), record the verdict in the \texttt{sampled} flag, and
propagate it so every downstream hop agrees. It is cheap, deterministic per
trace, and statistically unbiased --- and it throws away almost all of the
rare, interesting traces: at $p = 1\%$, you keep roughly 1\% of your
\emph{error} traces too. \textbf{Tail sampling} decides at the collector
after the whole trace is buffered: keep every trace containing an error or
a p99-plus latency, sample the healthy rest. It keeps exactly the traces
you debug with, at the price of collector memory, decision latency, and
the operational burden of buffering entire traces. The pragmatic hybrid:
head-sample at a few percent for unbiased statistics, tail-sample at the
collector to guarantee error and slow traces survive, and always record
counters \emph{before} sampling so metrics stay exact regardless.

\begin{warningbox}
Sampling happens \emph{after} metrics, never before. If your request
counters and latency histograms are derived from sampled spans, your RED
metrics are wrong by the sampling factor --- silently. Metrics must be
aggregated from every request; traces may be sampled.
\end{warningbox}
\subsection{Metrics Engineering: Histograms, Percentiles, Cardinality}
\label{subsec:ch18-metrics-engineering}

\paragraph{Why latency is a histogram, never an average.} API latency is
long-tailed: most requests are fast, a few are catastrophically slow, and
the slow ones are your most valuable users doing the largest operations.
The average hides exactly that. The fix is to record durations into
\emph{histo\-gram buckets} --- counters for $(0,5]$, $(5,10]$, $(10,25]$,
$(25,50]$, $(50,100]$, $(100,250]$, $(250,500]$, $(500,1000]$,
$(1000,2500]$, $(2500,\infty)$ milliseconds --- from which percentiles can
be estimated at query time.

\begin{definition}[Histogram percentile]
Given cumulative bucket counts $C_0 \le C_1 \le \dots \le C_k = N$ with
upper bounds $u_0 < u_1 < \dots < u_k$, the $q$-th percentile estimate is
the value interpolated within the first bucket where $C_j \ge
\lceil q N \rceil$:
\begin{equation}
\hat{p}_q \;=\; u_{j-1} + (u_j - u_{j-1})\,
\frac{\lceil qN \rceil - C_{j-1}}{C_j - C_{j-1}} .
\label{eq:ch18-histpct}
\end{equation}
Accuracy is bounded by bucket width: fine buckets where your SLO lives,
coarse buckets in the tail. Percentiles are \emph{estimated from counts},
never averaged from other percentiles.
\end{definition}

\begin{proposition}[The p99-of-averages fallacy, by counterexample]
\label{prop:ch18-p99avg}
Per-minute averages of a latency distribution can lie arbitrarily far below
the true p99. Take one minute with 1{,}000 requests: 999 at 10\,ms and one
at 10{,}000\,ms. The minute's average is
$\frac{999 \cdot 10 + 10000}{1000} \approx 20$\,ms. If every minute looks
like this, the ``p99 latency'' computed from per-minute averages is
$\approx 20$\,ms, while the true per-request p99 is $10{,}000$\,ms ---
a $500\times$ error, in the reassuring direction. Averaging first destroys
the tail; percentiles must be computed from per-request data (histograms),
and percentiles of percentiles or of averages are statistically void.
\end{proposition}

The same argument forbids averaging precomputed per-instance p99s across a
fleet: quantiles do not compose under averaging. Merging is legal only at
the histogram level (add the bucket counts, then apply
Equation~\eqref{eq:ch18-histpct}) --- which is precisely why the histogram
is the exportable unit.

\begin{figure}[ht]
  \centering
  \begin{tikzpicture}
    \begin{axis}[
        width=0.92\textwidth,
        height=7.2cm,
        ybar,
        bar width=10pt,
        xlabel={Latency bucket (ms)},
        ylabel={Requests},
        symbolic x coords={0--5, 5--10, 10--25, 25--50, 50--100,
                           100--250, 250--500, 500--1000, 1000--2500, 2500+},
        xtick=data,
        x tick label style={rotate=35, anchor=east, font=\footnotesize},
        ymin=0, ymax=4200,
        grid=major, grid style={gray!20},
        legend style={font=\footnotesize, at={(0.97,0.95)}, anchor=north east},
      ]
      \addplot[fill=darkblue!70, draw=darkblue] coordinates {
        (0--5, 320) (5--10, 1420) (10--25, 3980) (25--50, 2480)
        (50--100, 840) (100--250, 310) (250--500, 96)
        (500--1000, 210) (1000--2500, 190) (2500+, 154)
      };
      \addlegendentry{requests per bucket (seeded sample, $N=10{,}000$)}
      % mean ~761 ms, p50 ~25 ms, p99 ~23680 ms (Listing 18.4 demo output)
      \draw[thick, packtorange, dashed] (axis cs:50--100,0) -- (axis cs:50--100,4100)
        node[pos=0.97, anchor=south east, rotate=90, font=\footnotesize\bfseries, color=packtorange]
        {p50 $\approx$ 25 ms};
      \draw[thick, red!75!black, dashed] (axis cs:250--500,0) -- (axis cs:250--500,4100)
        node[pos=0.99, anchor=south east, rotate=90, font=\footnotesize\bfseries, color=red!75!black]
        {mean $\approx$ 761 ms};
      \draw[thick, purple, dashed] (axis cs:2500+,0) -- (axis cs:2500+,4100)
        node[pos=0.95, anchor=south west, rotate=90, font=\footnotesize\bfseries, color=purple]
        {p99 $\approx$ 23.7 s};
    \end{axis}
  \end{tikzpicture}
  \caption{The long tail, plotted from Listing~18.4's seeded distribution.
  Half of all requests finish in 25\,ms, yet the \emph{average} is
  761\,ms --- pulled 30$\times$ above the median by a tail the naked eye
  almost misses --- and p99 is three buckets past the average. Any
  dashboard or SLO keyed on the mean is measuring a user nobody has.}
  \label{fig:ch18-longtail}
\end{figure}

\paragraph{Cardinality discipline, with the cost math.} Every distinct
combination of label values is a new \emph{time series} that must be
stored, indexed, and scanned. Series count is multiplicative:

\begin{equation}
\text{series} \;=\; \prod_{\text{label } \ell} |\text{values}(\ell)| .
\label{eq:ch18-cardinality}
\end{equation}

A request counter labeled by route, status class, and region with
$5 \times 6 \times 3 = 90$ series is free. Add \texttt{user\_id} with
50{,}000 tenants and you have $90 \times 50{,}000 = 4.5$ million series;
at a few KB of index and chunk overhead per series, that is tens of GB of
memory and a query planner that hates you --- before counting ingestion
cost per active series per scrape interval. The rule: labels may only be
dimensions you would group by on a dashboard; anything unbounded
(user ids, request ids, full URLs, error messages) belongs in logs and
trace attributes, never in metric labels.

\begin{definition}[Exemplar]
An \emph{exemplar} is a trace id pinned onto a specific histogram bucket
observation. It is the bridge from ``p99 spiked'' (a metric) to ``here is
one concrete slow trace'' (a waterfall) --- the cheapest debugging
accelerator in the entire stack, and the reason metrics and traces must
share an ID space.
\end{definition}

\subsection{SLO Engineering: SLIs, Error Budgets, Burn Rates}
\label{subsec:ch18-slo}

\begin{definition}[SLI, SLO, error budget]
A \emph{service level indicator} (SLI) is a quantitative measure of service
behavior, expressed as a ratio of good events to total events. A
\emph{service level objective} (SLO) is a target on an SLI over a rolling
window (typically 30 days). The \emph{error budget} is the tolerated
shortfall: budget $= 1 - \text{SLO}$. For a 99.9\% availability SLO over a
30-day month, the budget is $0.001 \times 43{,}200\,\text{min} =
43.2$\,minutes of equivalent full outage per month~\cite{beyer2016sre}.
\end{definition}

Choosing SLIs for APIs is choosing a proxy for user pain. The three that
cover nearly every API:

\begin{itemize}
  \item \textbf{Availability:}
        $\text{SLI}_{\text{avail}} = \dfrac{\#\,\text{requests with status} < 500}
        {\#\,\text{requests}}$ per route. Count 5xx as failure; 4xx is the
        client's budget problem, not yours (Chapter~\ref{ch:problem_details}
        responses excluded).
  \item \textbf{Latency:}
        $\text{SLI}_{\text{lat}} = \dfrac{\#\,\text{requests faster than } T}
        {\#\,\text{requests}}$, computed from histogram buckets against the
        SLO threshold $T$ --- e.g.\ ``99\% of \texttt{GET /fleet/status}
        under 300\,ms.'' This is a threshold-count ratio, which is exactly
        what Equation~\eqref{eq:ch18-histpct}'s histogram answers.
  \item \textbf{Freshness} (data-serving APIs):
        $\text{SLI}_{\text{fresh}} = \dfrac{\#\,\text{reads served with data newer than } D}
        {\#\,\text{reads}}$ --- the right SLI when the API's job is
        recency (fleet positions, prices), where ``fast but stale'' is a
        silent failure mode.
\end{itemize}

\begin{table}[ht]
  \centering
  \small
  \begin{tabularx}{\textwidth}{l l X}
    \toprule
    \textbf{SLO} & \textbf{Budget / 30 days} & \textbf{What it permits} \\
    \midrule
    99\%    & 7 h 12 min  & A bad deploy \emph{and} a slow rollback, monthly. \\
    99.9\%  & 43.2 min    & One memorable incident; change freezes near zero. \\
    99.95\% & 21.6 min    & Pages must be rare and mitigation rehearsed. \\
    99.99\% & 4.32 min    & Automation mitigates; humans mostly supervise. \\
    \bottomrule
  \end{tabularx}
  \caption{Error budgets by SLO target. The budget is a \emph{spending
  allowance} shared by outages, risky deploys, and experiments --- when it
  is spent, the roadmap tilts to reliability~\cite{beyer2016sre}.}
  \label{tab:ch18-budgets}
\end{table}

\paragraph{Burn-rate alerting: the math.} Alerting on raw error counts
pages for trivia and sleeps through budget catastrophes. The SRE
Workbook's answer is to alert on \emph{burn rate} --- how fast the budget
is being consumed relative to its allotted pace~\cite{beyer2016sre}:

\begin{equation}
\text{burn rate} \;=\; \frac{\text{observed error rate}}{\text{budget rate}}
\;=\; \frac{e_w}{1 - \text{SLO}} ,
\qquad
\text{budget consumed} \;=\; \text{burn rate} \times \frac{w}{43{,}200\ \text{min}} .
\label{eq:ch18-burnrate}
\end{equation}

A burn rate of 1 exhausts the budget exactly as the period ends; 14.4
exhausts it in $720 / 14.4 = 50$ hours --- i.e.\ 2\% of the monthly budget
gone in one hour. Alerting on a single window forces a bad trade: a long
window is precise but slow, a short window is fast but noisy. The
\emph{multi-window multi-burn-rate} pattern fires only when \textbf{both}
a long window and a short window exceed the same burn-rate threshold:

\begin{equation}
\text{alert} \iff
\bigl(\text{burn}_{\text{long}}(t) \ge b\bigr)
\;\land\;
\bigl(\text{burn}_{\text{short}}(t) \ge b\bigr) .
\label{eq:ch18-mwmb}
\end{equation}

The long window certifies the burn is sustained; the short window
certifies it is happening \emph{now} (so the alert auto-resolves and does
not page on stale history). The canonical table:

\begin{table}[ht]
  \centering
  \small
  \begin{tabularx}{\textwidth}{l c c c X}
    \toprule
    \textbf{Policy} & \textbf{Burn} & \textbf{Long} & \textbf{Short} & \textbf{Meaning} \\
    \midrule
    Fast page & 14.4$\times$ & 1 h  & 5 m  & 2\% of budget in an hour: page now. \\
    Slow page & 6$\times$    & 6 h  & 30 m & 5\% of budget in six hours: page soon. \\
    Ticket    & 1$\times$    & 3 d  & 6 h  & 10\% of budget in three days: work-queue. \\
    \bottomrule
  \end{tabularx}
  \caption{Multi-window multi-burn-rate alerting for a 99.9\% SLO (SRE
  Workbook parameters)~\cite{beyer2016sre}. Listing~18.5 implements and
  tests exactly this table.}
  \label{tab:ch18-burnrate}
\end{table}

\begin{figure}[ht]
\centering
\begin{tikzpicture}[font=\small, >=stealth]
  % time axis
  \draw[->, thick] (0,0) -- (12.6,0) node[below left] {\footnotesize $t$};
  \foreach \x/\lab in {0/now$-6$h, 5/now$-1$h, 11/now$-5$m, 12.2/now}{
    \draw[gray!60] (\x,0.08) -- (\x,-0.08);
    \node[note, below] at (\x,-0.08) {\footnotesize \lab};
  }
  % long window bracket
  \draw[decorate, decoration={brace, amplitude=5pt}, thick, darkblue]
    (5,1.7) -- (12.2,1.7) node[midway, above=6pt, darkblue]
    {long window (1 h): sustained?};
  % short window bracket
  \draw[decorate, decoration={brace, amplitude=5pt}, thick, packtorange]
    (11,0.9) -- (12.2,0.9) node[midway, above=5pt, packtorange]
    {short window (5 m): happening now?};
  % error-rate trace
  \draw[very thick, red!70!black] (0,0.25) -- (8.5,0.25) -- (8.6,1.25)
    -- (11.6,1.25) -- (11.7,0.35) -- (12.2,0.35);
  \node[note, text=red!70!black, anchor=south west] at (8.7,1.3)
    {\footnotesize error rate spike};
  % threshold line
  \draw[dashed, thick, purple] (0,0.62) -- (12.2,0.62)
    node[pos=0.02, above, font=\footnotesize, color=purple]
    {$b \times$ budget rate};
  % AND gate annotation
  \node[draw, rounded corners, fill=yellow!15, align=center] at (4.2,3.3)
    {\footnotesize page $\iff$ burn$_{1\text{h}} \ge 14.4\ \land\ $burn$_{5\text{m}} \ge 14.4$};
  \draw[->, gray] (4.2,3.0) -- (7.8,1.95);
\end{tikzpicture}
\caption{Burn-rate alert logic. The long window (blue brace) filters noise
by requiring a sustained burn; the short window (orange brace) requires the
burn to be active \emph{at this instant}, so alerts clear promptly after
mitigation instead of paging on history. Both must hold
(Equation~\eqref{eq:ch18-mwmb}).}
\label{fig:ch18-burnrate}
\end{figure}

\paragraph{Alert fatigue is a system property, not a personality flaw.}
Every page that does not require immediate human action teaches the on-call
to ignore pages; enough false positives and a real page gets snoozed. The
discipline: page only on user-facing symptoms measured as budget burn
(fast/slow policies above); route everything else to tickets or dashboards;
delete or demote any alert that fired three times without action;
and review the page budget monthly like any other budget. An alert that
cannot name the SLO it protects is a dashboard panel wearing a pager
costume.

\subsection{Operational Practice: From Telemetry to Response}
\label{subsec:ch18-ops}

\paragraph{Correlation IDs end to end.} Chapter~\ref{ch:problem_details}
gave errors a \texttt{traceId}; this chapter gives it a source. The
discipline: accept or mint a \texttt{traceparent} at the edge, propagate it
through every internal call (headers for HTTP, message metadata for
queues), stamp it on every log line, echo it in every problem-details body
and support-facing error. The debugging loop becomes: client quotes id
$\to$ one log query returns the whole request's story $\to$ one trace
lookup returns the cross-service waterfall.

\paragraph{Dashboards per endpoint, anchored to SLOs.} A service dashboard
that cannot answer ``are we burning budget, and which endpoint is burning
it?'' is decoration. The layout that survives contact: row one is the SLO
(28-day SLI, budget remaining, burn-rate sparkline); row two is RED per
endpoint (rate, error ratio, duration percentiles from histograms); row
three is USE for the resources (pool utilization, queue depth); row four is
dependency health. Every panel cites the query; every query aggregates
only low-cardinality labels.

\paragraph{Debugging production latency systematically.} When p99 for
\texttt{POST /fleet/telemetry} jumps from 180\,ms to 900\,ms, resist the
log dive. Decompose the timeline with the trace waterfall: total span
(900\,ms) minus child spans ($\approx$ time in gateway code) $\to$ which
child dominates? If \texttt{orders.create} grew, recurse; if children are
unchanged but \emph{gaps} between them grew, the time went to queuing,
connection acquisition, TLS, or scheduling --- none of which appear as
spans unless you instrument them. The sequence is always: exemplar from
the violated bucket $\to$ waterfall $\to$ dominant span or gap $\to$
USE metrics for that span's resources $\to$ only then, correlated logs.

\paragraph{Incident response.} Severity is declared in user-impact units,
roles are assigned before diagnosis, and the postmortem is blameless or it
is fiction~\cite{beyer2016sre}. Northwind Robotics' table:

\begin{table}[ht]
  \centering
  \small
  \begin{tabularx}{\textwidth}{l X l}
    \toprule
    \textbf{Severity} & \textbf{Definition} & \textbf{Response} \\
    \midrule
    SEV1 & Total outage or data loss; SLO budget incinerating. & All-hands page; incident commander appointed in minutes. \\
    SEV2 & Major degradation of a core SLO; sustained fast burn. & Page owning team; commander + comms lead. \\
    SEV3 & Minor degradation or single-tenant impact. & Ticket; next business day. \\
    SEV4 & Cosmetic or internal-only. & Backlog. \\
    \bottomrule
  \end{tabularx}
  \caption{Severity levels. The response column is part of the definition:
  a severity without a response contract is a label, not a level.}
  \label{tab:ch18-severity}
\end{table}

The roles that matter at 3\,a.m.: the \emph{incident commander} owns
coordination and decisions (and does not debug); the \emph{operations lead}
debugs and mitigates; the \emph{communications lead} updates stakeholders
on a cadence. Runbooks --- per service, with SLOs, triage steps, rollback
commands, and escalation paths --- turn the first ten minutes from
improvisation into checklist. And the blameless postmortem: timeline,
mechanism-level root cause (``the deploy changed the connection-pool
default,'' never ``Alex broke the pool''), what went well, what went
poorly, action items with owners, and an explicit lesson for the SLO
itself --- did the alerts fire at the right time? Listing~18.6 generates
both skeletons from metadata, because structure enforced by code survives
stress better than structure enforced by memory.
%% ----------------------------------------------------------------------------
\section{Production-Ready Code Implementation}
\label{sec:ch18-code}

Everything below is typed Python~3.11, offline-first, and deterministic:
the OTLP exporter is replaced by a JSON-lines file exporter, metrics are
read back from an in-memory reader, trace ids come from a seeded generator,
and every timeline is seeded. The universe is Northwind Robotics fleet
telemetry. Save the listings under the indicated filenames and the printed
commands reproduce every claim in this chapter.

\subsection{Listing 18.1 --- The W3C Trace Context, by Hand}
\label{subsec:ch18-listing-traceparent}

\texttt{trace\_context.py} is a stdlib-only parser and injector for
\texttt{traceparent}, \texttt{tracestate}, and baggage. The design rule is
strictness: a malformed header raises \texttt{ValueError} naming the
offending field --- silently repairing a broken header would fork the
trace, and silently dropping it would hide a producer bug.

\noindent\textbf{Listing 18.1 --- \texttt{trace\_context.py}: W3C header
mechanics, stdlib only.}
\label{lst:ch18-trace-context}
\begin{minted}{python}
"""Hand-rolled W3C Trace Context: parse and inject ``traceparent``.

Stdlib-only reference implementation of the wire format every API engineer
should be able to read cold:

    traceparent: 00-4bf92f3577b34da6a3ce929d0e0e4736-00f067aa0ba902b7-01
                 |  |-------------------------------| |--------------| |-|
                 |  trace-id (32 hex, != 0)           parent/span-id   flags
                 version (2 hex, never "ff")          (16 hex, != 0)   bit 0 = sampled

Design notes:
  * Parsing is strict: malformed headers raise ``ValueError`` with the
    offending field named. Silently "fixing" a broken header would fork
    the trace; dropping it (returning ``None``) hides producer bugs.
  * Version "00" headers must be exactly 55 characters / 4 fields. Future
    versions may append fields, so for version != "00" we accept extra
    trailing fields and ignore them (forward compatibility).
  * ``tracestate`` and ``baggage`` parsers are included because the three
    headers travel together in production.

No network, no dependencies, no randomness: deterministic by construction.
"""

from __future__ import annotations

import secrets
from dataclasses import dataclass

HEX = frozenset("0123456789abcdef")
MAX_TRACESTATE_MEMBERS = 32  # per the W3C Trace Context spec


@dataclass(frozen=True)
class TraceParent:
    """An immutable, validated ``traceparent`` value."""

    version: str
    trace_id: str
    span_id: str
    sampled: bool

    def __post_init__(self) -> None:
        if len(self.version) != 2 or not set(self.version) <= HEX:
            raise ValueError(f"version must be 2 lowercase hex chars, got {self.version!r}")
        if self.version == "ff":
            raise ValueError("version 'ff' is reserved (invalid) by the spec")
        _check_hex_field("trace_id", self.trace_id, 32)
        _check_hex_field("span_id", self.span_id, 16)
        if self.trace_id == "0" * 32:
            raise ValueError("trace_id of all zeros is invalid")
        if self.span_id == "0" * 16:
            raise ValueError("span_id of all zeros is invalid")

    @property
    def flags(self) -> str:
        return "01" if self.sampled else "00"

    def header(self) -> str:
        """Render the canonical ``traceparent`` header value."""
        return f"{self.version}-{self.trace_id}-{self.span_id}-{self.flags}"


def _check_hex_field(name: str, value: str, length: int) -> None:
    if len(value) != length or not set(value) <= HEX:
        raise ValueError(
            f"{name} must be exactly {length} lowercase hex chars, got {value!r}"
        )


def parse_traceparent(header: str) -> TraceParent:
    """Parse a ``traceparent`` header value, raising ``ValueError`` if bad."""
    if not isinstance(header, str) or not header:
        raise ValueError("traceparent header must be a non-empty string")
    if header != header.strip():
        raise ValueError("traceparent must not carry surrounding whitespace")
    parts = header.split("-")
    if len(parts) < 4:
        raise ValueError(
            f"traceparent needs at least 4 dash-separated fields, got {len(parts)}"
        )
    version, trace_id, span_id, flags = parts[0], parts[1], parts[2], parts[3]
    if len(flags) != 2 or not set(flags) <= HEX:
        raise ValueError(f"flags must be 2 lowercase hex chars, got {flags!r}")
    if version == "00" and len(parts) != 4:
        raise ValueError("version '00' forbids extra fields")
    sampled = bool(int(flags, 16) & 0b0000_0001)  # only bit 0 is defined
    return TraceParent(version=version, trace_id=trace_id,
                       span_id=span_id, sampled=sampled)


def new_traceparent(*, sampled: bool = True) -> TraceParent:
    """Mint a fresh root context (16-byte trace id, 8-byte span id)."""
    return TraceParent(
        version="00",
        trace_id=secrets.token_hex(16),
        span_id=secrets.token_hex(8),
        sampled=sampled,
    )


def child_of(parent: TraceParent, *, sampled: bool | None = None) -> TraceParent:
    """Derive the outgoing context for a child hop: same trace, new span id."""
    return TraceParent(
        version=parent.version,
        trace_id=parent.trace_id,
        span_id=secrets.token_hex(8),
        sampled=parent.sampled if sampled is None else sampled,
    )


def parse_tracestate(header: str) -> dict[str, str]:
    """Parse a ``tracestate`` list-member header into an ordered dict.

    Raises ``ValueError`` on more than 32 members or malformed entries.
    """
    if not header.strip():
        return {}
    members = [m.strip() for m in header.split(",")]
    if len(members) > MAX_TRACESTATE_MEMBERS:
        raise ValueError(f"tracestate allows at most {MAX_TRACESTATE_MEMBERS} members")
    out: dict[str, str] = {}
    for member in members:
        if "=" not in member:
            raise ValueError(f"tracestate member lacks '=': {member!r}")
        key, _, value = member.partition("=")
        if not key or key in out:
            raise ValueError(f"tracestate duplicate/empty key: {key!r}")
        out[key] = value
    return out


def parse_baggage(header: str) -> dict[str, str]:
    """Parse a W3C ``baggage`` header; unknown per-entry properties dropped."""
    out: dict[str, str] = {}
    for member in header.split(","):
        member = member.strip()
        if not member:
            continue
        pair = member.split(";")[0].strip()  # drop ;property annotations
        if "=" not in pair:
            raise ValueError(f"baggage member lacks '=': {member!r}")
        key, _, value = pair.partition("=")
        out[key.strip()] = value.strip()
    return out


def _demo() -> None:
    """Round-trip one trace through three synthetic hops (deterministic)."""
    print("== W3C Trace Context round-trip demo (Northwind Robotics) ==")
    root = parse_traceparent(
        "00-4bf92f3577b34da6a3ce929d0e0e4736-00f067aa0ba902b7-01"
    )
    print(f"incoming  : {root.header()}  sampled={root.sampled}")
    hop1 = child_of(root)
    print(f"gateway -> orders      : {hop1.header()}")
    hop2 = child_of(hop1)
    print(f"orders  -> inventory   : {hop2.header()}")
    echoed = parse_traceparent(hop2.header())
    assert echoed.trace_id == root.trace_id, "trace id must survive every hop"
    assert echoed.span_id == hop2.span_id, "round-trip must preserve span id"
    print("trace id stable across hops:", echoed.trace_id == root.trace_id)
    print(parse_tracestate("nr=region:eu-1,congo=t61rcWkgMzE"))
    print(parse_baggage("tenant=northwind-robotics,plan=fleet-pro;property=x"))


if __name__ == "__main__":
    _demo()
\end{minted}

\subsection{Listing 18.2 --- The Header Mechanics, Proven}
\label{subsec:ch18-listing-traceparent-tests}

Round-trip, rejection, forward-compatibility, and propagation tests --- the
wire format as an executable specification.

\noindent\textbf{Listing 18.2 --- \texttt{test\_trace\_context.py}: 19
tests for the trace-context contract.}
\label{lst:ch18-trace-context-tests}
\begin{minted}{python}
"""Executable proof sheet for trace_context.py (run: pytest -q)."""

from __future__ import annotations

import pytest

from trace_context import (
    TraceParent,
    child_of,
    new_traceparent,
    parse_baggage,
    parse_traceparent,
    parse_tracestate,
)

VALID = "00-4bf92f3577b34da6a3ce929d0e0e4736-00f067aa0ba902b7-01"


def test_round_trip_preserves_every_field() -> None:
    tp = parse_traceparent(VALID)
    assert tp.header() == VALID
    assert tp.version == "00"
    assert tp.trace_id == "4bf92f3577b34da6a3ce929d0e0e4736"
    assert tp.span_id == "00f067aa0ba902b7"
    assert tp.sampled is True


def test_not_sampled_flag() -> None:
    tp = parse_traceparent(VALID[:-2] + "00")
    assert tp.sampled is False
    assert tp.header().endswith("-00")


def test_flags_bitmask_ignores_undefined_bits() -> None:
    # only bit 0 (0x01) is the sampled flag; other bits must not leak in
    tp = parse_traceparent(VALID[:-2] + "03")
    assert tp.sampled is True
    assert tp.flags == "01"  # canonical rendering keeps only defined bits


def test_child_keeps_trace_id_and_gets_fresh_span_id() -> None:
    root = parse_traceparent(VALID)
    child = child_of(root)
    assert child.trace_id == root.trace_id
    assert child.span_id != root.span_id
    assert len(child.span_id) == 16
    assert child.sampled == root.sampled


def test_new_contexts_are_well_formed() -> None:
    tp = new_traceparent()
    again = parse_traceparent(tp.header())
    assert again == tp


@pytest.mark.parametrize(
    "bad",
    [
        "",                                   # empty
        VALID.upper(),                        # uppercase hex forbidden
        "ff-4bf92f3577b34da6a3ce929d0e0e4736-00f067aa0ba902b7-01",  # ff version
        "00-00000000000000000000000000000000-00f067aa0ba902b7-01",  # zero trace id
        "00-4bf92f3577b34da6a3ce929d0e0e4736-0000000000000000-01",  # zero span id
        "00-4bf92f3577b34da6a3ce929d0e0e473-00f067aa0ba902b7-01",   # short trace id
        "00-4bf92f3577b34da6a3ce929d0e0e4736-00f067aa0ba902b7",     # missing flags
        "00-4bf92f3577b34da6a3ce929d0e0e4736-00f067aa0ba902b7-0g",  # bad hex flag
        "00-4bf92f3577b34da6a3ce929d0e0e4736-00f067aa0ba902b7-01-extra",  # v00 extra field
        " 00-4bf92f3577b34da6a3ce929d0e0e4736-00f067aa0ba902b7-01",  # whitespace
    ],
)
def test_rejects_malformed_headers(bad: str) -> None:
    with pytest.raises(ValueError):
        parse_traceparent(bad)


def test_future_version_may_append_fields() -> None:
    # forward compatibility: version 01 with a trailing field parses fine
    tp = parse_traceparent("01-" + VALID[3:] + "-congo")
    assert tp.version == "01"


def test_tracestate_round_trip_and_limits() -> None:
    state = parse_tracestate("nr=region:eu-1,congo=t61rcWkgMzE")
    assert state == {"nr": "region:eu-1", "congo": "t61rcWkgMzE"}
    with pytest.raises(ValueError):
        parse_tracestate(",".join(f"k{i}=v" for i in range(33)))
    with pytest.raises(ValueError):
        parse_tracestate("nr=one,nr=two")  # duplicate key


def test_baggage_drops_properties() -> None:
    bag = parse_baggage("tenant=northwind-robotics;ttl=30,plan=fleet-pro")
    assert bag == {"tenant": "northwind-robotics", "plan": "fleet-pro"}
    with pytest.raises(ValueError):
        parse_baggage("no-equals-sign")


def test_immutability() -> None:
    tp = parse_traceparent(VALID)
    with pytest.raises(AttributeError):
        tp.trace_id = "0" * 32  # type: ignore[misc]
\end{minted}

\subsection{Listing 18.3 --- An Instrumented FastAPI Service, Offline}
\label{subsec:ch18-listing-app}

\texttt{instrumented\_app.py} wires all three pillars into one service. A
seeded ID generator makes span ids reproducible; \texttt{JsonlSpanExporter}
implements the real \texttt{SpanExporter} interface but writes spans to a
local file (the offline OTLP replacement); an \texttt{InMemoryMetricReader}
holds RED metrics (counter plus explicit-bucket latency histogram); and a
logging filter stamps the active span's ids onto every record. The
three-hop chain --- gateway $\to$ orders $\to$ inventory --- propagates
context by injecting and extracting real \texttt{traceparent} strings
between hops, and the demo \emph{asserts} the exported span tree: one trace
id, correct parent chain, correct span kinds. The FastAPI app itself is
fully wired (instrumentor plus correlation middleware); the demo drives the
service layer directly so the whole thing runs with no server and no
network.

\noindent\textbf{Listing 18.3 --- \texttt{instrumented\_app.py}: traces,
metrics, and correlated logs with propagation proof.}
\label{lst:ch18-instrumented-app}
\begin{minted}{python}
"""Instrumented FastAPI service: traces, metrics, and correlated logs.

Northwind Robotics fleet-telemetry platform, offline edition. The OTLP
exporter you would point at a real collector is replaced by
``JsonlSpanExporter`` (spans to a local JSON-lines file) and an
``InMemoryMetricReader`` (metrics read back from the SDK), so the whole
demo runs with no collector, no network, and no credentials.

Telemetry wiring:
  * one ``TracerProvider`` with a seeded ID generator (deterministic tests)
  * one ``MeterProvider`` with explicit-bucket histogram views (RED metrics)
  * a ``logging.Formatter`` that stamps every record with trace_id/span_id
    of the *active* span, so a log line and a trace never drift apart

Propagation demo: a synthetic incoming ``traceparent`` header is extracted
at the gateway, and each internal hop injects the current context into a
fresh carrier dict --- the exact bytes that would ride the wire in
production. The exported span tree proves all three hops share one
trace_id and nest parent -> child correctly.

Run:  python instrumented_app.py
Serve (optional, needs uvicorn):  uvicorn instrumented_app:app
"""

from __future__ import annotations

import json
import logging
import random
import time
from dataclasses import dataclass, field
from pathlib import Path
from typing import Mapping, Sequence

from fastapi import FastAPI, Request
from fastapi.responses import JSONResponse
from opentelemetry import trace
from opentelemetry.instrumentation.fastapi import FastAPIInstrumentor
from opentelemetry.sdk.metrics import MeterProvider
from opentelemetry.sdk.metrics.export import InMemoryMetricReader
from opentelemetry.sdk.metrics.view import ExplicitBucketHistogramAggregation, View
from opentelemetry.sdk.resources import Resource
from opentelemetry.sdk.trace import TracerProvider
from opentelemetry.sdk.trace.export import (
    SimpleSpanProcessor,
    SpanExporter,
    SpanExportResult,
)
from opentelemetry.sdk.trace.id_generator import IdGenerator
from opentelemetry.sdk.trace.sampling import ALWAYS_ON
from opentelemetry.trace.propagation.tracecontext import (
    TraceContextTextMapPropagator,
)
from opentelemetry.trace import SpanKind

LATENCY_BUCKETS_MS = [5.0, 10.0, 25.0, 50.0, 100.0, 250.0, 500.0, 1000.0, 2500.0]
PROPAGATOR = TraceContextTextMapPropagator()


# --------------------------------------------------------------------------
# Telemetry plumbing
# --------------------------------------------------------------------------
class SeededIdGenerator(IdGenerator):
    """Deterministic trace/span IDs so tests assert exact span trees."""

    def __init__(self, seed: int) -> None:
        self._rng = random.Random(seed)

    def generate_trace_id(self) -> int:
        return max(1, self._rng.getrandbits(128))  # all-zero id is invalid

    def generate_span_id(self) -> int:
        return max(1, self._rng.getrandbits(64))


class JsonlSpanExporter(SpanExporter):
    """Offline OTLP replacement: append finished spans to a JSONL file."""

    def __init__(self, path: Path) -> None:
        self._path = path
        self._path.write_text("", encoding="utf-8")

    def export(self, spans: Sequence[trace.Span]) -> SpanExportResult:
        with self._path.open("a", encoding="utf-8") as fh:
            for span in spans:
                fh.write(json.dumps(_span_to_dict(span)) + "\n")
        return SpanExportResult.SUCCESS

    def shutdown(self) -> None:  # nothing buffered; file handles are scoped
        return None


def _span_to_dict(span: trace.Span) -> dict[str, object]:
    ctx = span.get_span_context()
    parent = span.parent
    return {
        "name": span.name,
        "kind": span.kind.name if span.kind else "INTERNAL",
        "trace_id": f"{ctx.trace_id:032x}",
        "span_id": f"{ctx.span_id:016x}",
        "parent_span_id": f"{parent.span_id:016x}" if parent else None,
        "duration_ms": round((span.end_time - span.start_time) / 1e6, 3),
        "attributes": {k: v for k, v in (span.attributes or {}).items()},
        "status": span.status.status_code.name,
    }


class TraceContextFilter(logging.Filter):
    """Stamps trace_id/span_id of the active span onto every record."""

    def filter(self, record: logging.LogRecord) -> bool:
        ctx = trace.get_current_span().get_span_context()
        record.trace_id = f"{ctx.trace_id:032x}" if ctx.trace_id else "-"
        record.span_id = f"{ctx.span_id:016x}" if ctx.span_id else "-"
        return True


class JsonFormatter(logging.Formatter):
    """One JSON object per line; never log headers, tokens, or payloads."""

    def format(self, record: logging.LogRecord) -> str:
        return json.dumps({
            "ts": self.formatTime(record, "%Y-%m-%dT%H:%M:%S"),
            "level": record.levelname,
            "logger": record.name,
            "msg": record.getMessage(),
            "trace_id": getattr(record, "trace_id", "-"),
            "span_id": getattr(record, "span_id", "-"),
        })


@dataclass
class Telemetry:
    """Everything an instrumented service needs, wired offline."""

    provider: TracerProvider
    tracer: trace.Tracer
    exporter_path: Path
    meter_reader: InMemoryMetricReader
    request_counter: object
    request_duration: object
    logger: logging.Logger
    log_lines: list[str] = field(default_factory=list)


def build_telemetry(seed: int = 18, spans_path: Path = Path("spans.jsonl")) -> Telemetry:
    resource = Resource.create({"service.name": "nr-fleet-telemetry"})
    provider = TracerProvider(
        sampler=ALWAYS_ON,
        id_generator=SeededIdGenerator(seed),
        resource=resource,
    )
    provider.add_span_processor(SimpleSpanProcessor(JsonlSpanExporter(spans_path)))

    reader = InMemoryMetricReader()
    view = View(
        instrument_name="http.server.request.duration",
        aggregation=ExplicitBucketHistogramAggregation(boundaries=LATENCY_BUCKETS_MS),
    )
    meter_provider = MeterProvider(
        metric_readers=[reader], views=[view], resource=resource
    )
    meter = meter_provider.get_meter("nr.fleet")

    logger = logging.getLogger("nr.fleet")
    logger.setLevel(logging.INFO)
    logger.handlers.clear()
    logger.propagate = False

    telemetry = Telemetry(
        provider=provider,
        tracer=provider.get_tracer("nr.fleet"),
        exporter_path=spans_path,
        meter_reader=reader,
        request_counter=meter.create_counter(
            "http.server.request.count",
            description="Total HTTP requests (RED: rate, errors)",
        ),
        request_duration=meter.create_histogram(
            "http.server.request.duration",
            unit="ms",
            description="Request latency histogram (RED: duration)",
        ),
        logger=logger,
    )
    return telemetry


# --------------------------------------------------------------------------
# Service layer: a three-hop internal call chain
# --------------------------------------------------------------------------
def inventory_reserve(telemetry: Telemetry, carrier: dict[str, str]) -> dict[str, object]:
    """Hop 3 (deepest): check stock. Extracts context from its carrier."""
    ctx = PROPAGATOR.extract(carrier=carrier)
    with telemetry.tracer.start_as_current_span(
        "inventory.reserve", context=ctx, kind=SpanKind.SERVER
    ) as span:
        span.set_attribute("peer.service", "inventory")
        span.set_attribute("warehouse", "eu-1")
        time.sleep(0.004)  # deterministic-ish work stand-in; tests don't time it
        telemetry.logger.info("stock reserved sku=NR-GRIP-4 qty=2")
        return {"sku": "NR-GRIP-4", "reserved": 2}


def orders_create(telemetry: Telemetry, carrier: dict[str, str]) -> dict[str, object]:
    """Hop 2: create the resupply order; injects context for hop 3."""
    ctx = PROPAGATOR.extract(carrier=carrier)
    with telemetry.tracer.start_as_current_span(
        "orders.create", context=ctx, kind=SpanKind.SERVER
    ) as span:
        span.set_attribute("peer.service", "orders")
        span.set_attribute("order.kind", "resupply")
        outbound: dict[str, str] = {}
        PROPAGATOR.inject(outbound)  # traceparent bytes for the next hop
        reservation = inventory_reserve(telemetry, outbound)
        telemetry.logger.info("order created id=NR-2026-0917")
        return {"order_id": "NR-2026-0917", "reservation": reservation}


def gateway_handle(telemetry: Telemetry, headers: Mapping[str, str]) -> dict[str, object]:
    """Hop 1 (edge): extract the client-supplied traceparent, record RED."""
    started = time.perf_counter()
    status = 200
    ctx = PROPAGATOR.extract(carrier=dict(headers))
    with telemetry.tracer.start_as_current_span(
        "gateway POST /fleet/telemetry", context=ctx, kind=SpanKind.SERVER
    ) as span:
        span.set_attribute("http.request.method", "POST")
        span.set_attribute("url.path", "/fleet/telemetry")
        span.set_attribute("http.response.status_code", status)
        outbound: dict[str, str] = {}
        PROPAGATOR.inject(outbound)
        try:
            result = orders_create(telemetry, outbound)
        except Exception:
            status = 500
            span.set_attribute("http.response.status_code", status)
            raise
        finally:
            elapsed_ms = (time.perf_counter() - started) * 1000.0
            attrs = {"route": "/fleet/telemetry", "status": status}
            telemetry.request_counter.add(1, attrs)
            telemetry.request_duration.record(elapsed_ms, attrs)
        telemetry.logger.info("request complete status=200")
        return result


# --------------------------------------------------------------------------
# FastAPI wiring (thin adapters over the instrumented service layer)
# --------------------------------------------------------------------------
telemetry = build_telemetry()
app = FastAPI(title="Northwind Robotics Fleet Telemetry")
FastAPIInstrumentor.instrument_app(
    app, tracer_provider=telemetry.provider, excluded_urls="/health"
)


@app.middleware("http")
async def correlate_logs(request: Request, call_next):  # noqa: ANN001
    # The instrumentor's server span is active here, so logs pick up its ids.
    response: JSONResponse = await call_next(request)
    telemetry.logger.info(
        "http request method=%s path=%s status=%s",
        request.method, request.url.path, response.status_code,
    )
    return response


@app.post("/fleet/telemetry")
async def ingest(request: Request) -> dict[str, object]:
    return gateway_handle(telemetry, dict(request.headers))


@app.get("/health")
async def health() -> dict[str, str]:
    return {"status": "ok"}


# --------------------------------------------------------------------------
# Offline demo: drive the service layer with a synthetic traceparent
# --------------------------------------------------------------------------
def _demo() -> None:
    telemetry = build_telemetry(seed=18, spans_path=Path("spans.jsonl"))

    handler = logging.StreamHandler()
    handler.setFormatter(JsonFormatter())
    handler.addFilter(TraceContextFilter())
    telemetry.logger.addHandler(handler)

    incoming = {
        "traceparent": "00-4bf92f3577b34da6a3ce929d0e0e4736-00f067aa0ba902b7-01",
        "tracestate": "nr=region:eu-1",
    }
    print("== three-hop trace, one synthetic incoming traceparent ==")
    result = gateway_handle(telemetry, incoming)
    telemetry.provider.force_flush()

    lines = telemetry.exporter_path.read_text(encoding="utf-8").splitlines()
    spans = [json.loads(line) for line in lines if line.strip()]
    print(f"exported {len(spans)} spans:")
    for s in sorted(spans, key=lambda s: s["parent_span_id"] is not None):
        print(
            f"  {s['name']:<32} kind={s['kind']:<8} "
            f"trace={s['trace_id'][:12]}... span={s['span_id']} "
            f"parent={s['parent_span_id']}"
        )
    trace_ids = {s["trace_id"] for s in spans}
    by_name = {s["name"]: s for s in spans}
    assert len(trace_ids) == 1, "all hops must share one trace_id"
    assert trace_ids == {"4bf92f3577b34da6a3ce929d0e0e4736"}, (
        "incoming trace_id must survive every hop"
    )
    assert by_name["orders.create"]["parent_span_id"] == by_name[
        "gateway POST /fleet/telemetry"]["span_id"]
    assert by_name["inventory.reserve"]["parent_span_id"] == by_name[
        "orders.create"]["span_id"]
    print("propagation verified: one trace_id, correct parent chain")

    data = telemetry.meter_reader.get_metrics_data()
    for rm in data.resource_metrics:
        for sm in rm.scope_metrics:
            for m in sm.metrics:
                pts = list(m.data.data_points)
                print(f"metric {m.name}: points={len(pts)}")
    print(f"result: {json.dumps(result)}")


if __name__ == "__main__":
    _demo()
\end{minted}

\noindent Running it prints the proof of propagation (ids are seeded and
reproduce exactly):

\begin{minted}{text}
exported 3 spans:
  inventory.reserve                kind=SERVER   trace=4bf92f3577b3... parent=72e63ac7a9538322
  orders.create                    kind=SERVER   trace=4bf92f3577b3... parent=1f70d5dc2e675fc7
  gateway POST /fleet/telemetry    kind=SERVER   trace=4bf92f3577b3... parent=00f067aa0ba902b7
propagation verified: one trace_id, correct parent chain
\end{minted}

Note the gateway's parent: \texttt{00f067aa0ba902b7} --- the span id from
the synthetic incoming \texttt{traceparent} of Listing~18.1's example. The
client's trace context survived the whole chain.

\subsection{Listing 18.4 --- Histograms: Percentiles That Survive the Tail}
\label{subsec:ch18-listing-histogram}

\texttt{latency\_histograms.py} re-implements the HdrHistogram layout in
stdlib: power-of-two buckets with linear sub-buckets give bounded relative
error (two significant figures $\Rightarrow$ bucket width under 1\% of the
value) at $O(1)$ memory per recording. The embedded pytest suite proves the
two properties that matter: percentiles match exact rank-based quantiles
where resolution permits, and larger values stay within the resolution
bound. The demo prints the long-tail comparison behind
Figure~\ref{fig:ch18-longtail}: mean 761\,ms, p50 25\,ms, p99 23{,}680\,ms
--- the average lying about a 31$\times$ gap.

\noindent\textbf{Listing 18.4 --- \texttt{latency\_histograms.py}: log-bucket
histogram with embedded proof suite.}
\label{lst:ch18-histograms}
\begin{minted}{python}
"""Log-bucket latency histograms: percentiles that survive the long tail.

A stdlib-only re-implementation of the HdrHistogram layout (power-of-two
buckets with linear sub-buckets), because "compute p99" is a sentence every
API engineer says and few can defend. Properties:

  * O(1) memory per recording, fixed-size counts array
  * bounded relative error: 2 significant figures => every recorded value
    lands in a bucket whose width is at most ~0.8% of the value
  * percentiles are read from cumulative bucket counts, never from a
    sample list, so cost is O(buckets), not O(samples)

The demo generates a seeded long-tail distribution (typical API latency:
fast median, brutal tail) and prints the comparison that motivates this
whole chapter: the naive average sits comfortably below p99 --- the average
is lying to you about your users' pain.

Run:  python latency_histograms.py        (demo)
      pytest latency_histograms.py -q     (embedded test suite)
"""

from __future__ import annotations

import math
import random
from dataclasses import dataclass, field


@dataclass
class LogHistogram:
    """HdrHistogram-style integer histogram (values in milliseconds)."""

    lowest: int = 1
    highest: int = 3_600_000  # one hour in ms
    significant_figures: int = 2
    total: int = 0
    sum_ms: float = 0.0
    counts: list[int] = field(init=False)

    def __post_init__(self) -> None:
        if not 1 <= self.significant_figures <= 5:
            raise ValueError("significant_figures must be in 1..5")
        if self.lowest < 1 or self.highest < 2 * self.lowest:
            raise ValueError("need 1 <= lowest <= highest/2")
        # sub_bucket_count: smallest power of two >= 2 * 10^sig
        self._sub_mag = math.ceil(math.log2(2 * 10 ** self.significant_figures))
        self._sub_count = 1 << self._sub_mag
        self._half = self._sub_count // 2
        self._unit_mag = (self.lowest).bit_length() - 1
        # bucket_count: enough doublings to reach `highest`
        self._bucket_count = 1
        trackable = self._sub_count << self._unit_mag
        while trackable < self.highest:
            trackable <<= 1
            self._bucket_count += 1
        self.counts = [0] * ((self._bucket_count + 1) * self._half)

    # -- recording ---------------------------------------------------------
    def _counts_index(self, value: int) -> int:
        if value < 0 or value > self.highest:
            raise ValueError(f"value {value} outside [0, {self.highest}]")
        v = max(value >> self._unit_mag, 1)
        bucket = max(0, v.bit_length() - self._sub_mag)
        sub_bucket = v >> bucket
        return (bucket + 1) * self._half + (sub_bucket - self._half)

    def record(self, value: int) -> None:
        self.counts[self._counts_index(value)] += 1
        self.total += 1
        self.sum_ms += value

    # -- reading -----------------------------------------------------------
    def _value_at(self, counts_index: int) -> int:
        bucket = (counts_index >> (self._sub_mag - 1)) - 1
        sub_bucket = (counts_index & (self._half - 1)) + self._half
        if bucket < 0:
            bucket = 0
            sub_bucket -= self._half
        return sub_bucket << (bucket + self._unit_mag)

    def percentile(self, q: float) -> int:
        """Smallest recorded-bucket value at which q% of samples fall."""
        if not 0 < q <= 100:
            raise ValueError("q must be in (0, 100]")
        if self.total == 0:
            raise ValueError("no samples recorded")
        rank = max(1, math.ceil(q / 100.0 * self.total))
        cumulative = 0
        for idx, count in enumerate(self.counts):
            cumulative += count
            if cumulative >= rank:
                return self._value_at(idx)
        raise AssertionError("unreachable: rank <= total")

    def mean(self) -> float:
        if self.total == 0:
            raise ValueError("no samples recorded")
        return self.sum_ms / self.total


def exact_quantile(data: list[int], q: float) -> int:
    """Ground truth on small data: the rank-based quantile of a sorted list."""
    if not data:
        raise ValueError("empty data")
    ordered = sorted(data)
    rank = max(1, math.ceil(q / 100.0 * len(ordered)))
    return ordered[rank - 1]


def synthetic_latencies(seed: int = 18, n: int = 10_000) -> list[int]:
    """Seeded long-tail: 95% of requests are healthy, 5% hit the tail."""
    rng = random.Random(seed)
    out: list[int] = []
    for _ in range(n):
        if rng.random() < 0.95:  # healthy path: tight lognormal around ~25 ms
            out.append(max(1, int(rng.lognormvariate(math.log(25), 0.35))))
        else:                    # tail path: GC pauses, retries, cold caches
            out.append(rng.randint(800, 30_000))
    return out


def _demo() -> None:
    data = synthetic_latencies()
    hist = LogHistogram()
    for ms in data:
        hist.record(ms)
    print(f"samples={hist.total}  (seeded, long-tail)")
    print(f"  naive average : {hist.mean():9.1f} ms   <-- the lie")
    for q in (50, 95, 99, 99.9):
        print(f"  p{q:<5}        : {hist.percentile(q):9} ms")
    ratio = hist.percentile(99) / hist.mean()
    print(f"  p99 / average : {ratio:9.1f}x  <-- what the average was hiding")


# ---------------------------------------------------------------------------
# Embedded test suite: pytest latency_histograms.py -q
# ---------------------------------------------------------------------------
def test_percentiles_match_exact_quantiles_on_small_data() -> None:
    # values below sub_bucket_count (256) each get their own bucket,
    # so histogram percentiles must equal exact rank-based quantiles
    data = list(range(1, 256))
    hist = LogHistogram()
    for v in data:
        hist.record(v)
    for q in (1, 25, 50, 75, 90, 95, 99, 99.9, 100):
        assert hist.percentile(q) == exact_quantile(data, q)


def test_larger_values_stay_within_resolution_bound() -> None:
    data = list(range(1, 5001))
    hist = LogHistogram()
    for v in data:
        hist.record(v)
    for q in (25, 50, 75, 90, 99):
        approx, exact = hist.percentile(q), exact_quantile(data, q)
        assert abs(approx - exact) / exact < 0.01  # 2 significant figures


def test_mean_is_exact() -> None:
    hist = LogHistogram()
    for v in [10, 20, 30]:
        hist.record(v)
    assert hist.mean() == 20.0


def test_rank_rounding_uses_ceiling() -> None:
    data = [5] * 9 + [900]  # p90 rank = ceil(9.0) = 9 -> 5; p95 -> 900
    hist = LogHistogram()
    for v in data:
        hist.record(v)
    assert hist.percentile(90) == 5
    assert hist.percentile(95) == 900


def test_relative_error_is_bounded_by_resolution() -> None:
    hist = LogHistogram(significant_figures=2)
    hist.record(12_345)  # bucket width at this magnitude is 128 ms
    approx = hist.percentile(50)
    assert abs(approx - 12_345) / 12_345 < 0.01


def test_out_of_range_and_empty_rejected() -> None:
    hist = LogHistogram()
    import pytest

    with pytest.raises(ValueError):
        hist.record(hist.highest + 1)
    with pytest.raises(ValueError):
        hist.percentile(99)
    with pytest.raises(ValueError):
        hist.percentile(0)


def test_long_tail_demo_is_deterministic() -> None:
    assert synthetic_latencies() == synthetic_latencies()
    hist = LogHistogram()
    for ms in synthetic_latencies():
        hist.record(ms)
    assert hist.percentile(99) > 10 * hist.mean()  # the tail really is long


if __name__ == "__main__":
    _demo()
\end{minted}

\subsection{Listing 18.5 --- The Burn-Rate Simulator}
\label{subsec:ch18-listing-burnrate}

\texttt{burn\_rate\_sim.py} implements Equations~\eqref{eq:ch18-burnrate}
and~\eqref{eq:ch18-mwmb} against a seeded three-day timeline with a
scripted incident (50 minutes at an 80$\times$ burn --- a wedged
dependency). Window queries run over prefix sums, so evaluation is $O(n)$.
The demo prints the policy table and the alert timeline; the tests prove
the budget arithmetic (43.2 minutes; 2\%/5\%/10\% consumption), that
healthy traffic never pages, that the incident pages fast then slow but
never tickets, and that a sustained 3$\times$ burn tickets without paging.

\noindent\textbf{Listing 18.5 --- \texttt{burn\_rate\_sim.py}: multi-window
multi-burn-rate alerting, simulated offline.}
\label{lst:ch18-burnrate}
\begin{minted}{python}
"""Multi-window, multi-burn-rate SLO alerting simulator (offline, seeded).

Feeds a synthetic minute-by-minute traffic timeline for the Northwind
Robotics fleet API into the alerting logic from the Google SRE Workbook's
"Alerting on SLOs" chapter: an alert fires only when BOTH a long window
and a short window show a burn rate above the policy threshold.

    burn_rate(t, w)   = error_rate(t, w) / budget_rate
    budget_rate       = 1 - SLO                (0.001 for a 99.9% SLO)
    budget_consumed   = burn_rate * w / period (% of the 30-day budget)

Policy table (budget consumed if the burn persists for the long window):
    fast page  : burn 14.4  long 1h   short 5m   -> 2%  of budget
    slow page  : burn  6.0  long 6h   short 30m  -> 5%  of budget
    ticket     : burn  1.0  long 3d   short 6h   -> 10% of budget

Run:  python burn_rate_sim.py        (scripted incident demo)
      pytest burn_rate_sim.py -q     (embedded test suite)
"""

from __future__ import annotations

import random
from dataclasses import dataclass

SLO = 0.999
BUDGET_RATE = 1.0 - SLO              # 0.001
PERIOD_MINUTES = 30 * 24 * 60        # 30-day rolling window: 43,200 min
BUDGET_MINUTES_EQUIVALENT = BUDGET_RATE * PERIOD_MINUTES  # 43.2 min


@dataclass(frozen=True)
class MinuteSample:
    minute: int
    requests: int
    errors: int


@dataclass(frozen=True)
class BurnRatePolicy:
    name: str
    burn_rate: float
    long_window_min: int
    short_window_min: int

    @property
    def budget_consumed_pct(self) -> float:
        """% of the period budget burned if this rate lasts the long window."""
        return 100.0 * self.burn_rate * self.long_window_min / PERIOD_MINUTES


POLICIES = [
    BurnRatePolicy("fast-page", 14.4, 60, 5),
    BurnRatePolicy("slow-page", 6.0, 360, 30),
    BurnRatePolicy("ticket", 1.0, 3 * 24 * 60, 360),
]


@dataclass(frozen=True)
class Alert:
    policy: str
    fired_at: int
    long_burn: float
    short_burn: float


def window_rate(samples: list[MinuteSample], t: int, window: int) -> float:
    """Error rate over samples in (t - window, t]; 0.0 when no traffic."""
    total = errors = 0
    for s in samples:
        if t - window < s.minute <= t:
            total += s.requests
            errors += s.errors
    return errors / total if total else 0.0


def burn_rate(samples: list[MinuteSample], t: int, window: int) -> float:
    return window_rate(samples, t, window) / BUDGET_RATE


def _prefix_sums(samples: list[MinuteSample]) -> tuple[list[int], list[int]]:
    """Cumulative requests/errors per minute for O(1) window queries."""
    horizon = samples[-1].minute
    cum_req = [0] * (horizon + 1)
    cum_err = [0] * (horizon + 1)
    for s in samples:
        cum_req[s.minute] += s.requests
        cum_err[s.minute] += s.errors
    for t in range(1, horizon + 1):
        cum_req[t] += cum_req[t - 1]
        cum_err[t] += cum_err[t - 1]
    return cum_req, cum_err


def _burn(cum_req: list[int], cum_err: list[int], t: int, window: int) -> float:
    lo = max(0, t - window)  # window (t - window, t]
    total = cum_req[t] - cum_req[lo]
    errors = cum_err[t] - cum_err[lo]
    return (errors / total / BUDGET_RATE) if total else 0.0


def evaluate(
    samples: list[MinuteSample],
    policies: list[BurnRatePolicy] | None = None,
) -> list[Alert]:
    """First firing of each policy; reset requires the long window to clear."""
    fired: list[Alert] = []
    active: set[str] = set()
    if not samples:
        return fired
    policies = policies or POLICIES
    cum_req, cum_err = _prefix_sums(samples)
    for t in range(1, samples[-1].minute + 1):
        for policy in policies:
            if t < policy.long_window_min:
                continue  # window not yet fully populated; don't alert on cold data
            long_burn = _burn(cum_req, cum_err, t, policy.long_window_min)
            short_burn = _burn(cum_req, cum_err, t, policy.short_window_min)
            breached = (
                long_burn >= policy.burn_rate
                and short_burn >= policy.burn_rate
            )
            if breached and policy.name not in active:
                active.add(policy.name)
                fired.append(Alert(policy.name, t, long_burn, short_burn))
            elif long_burn < policy.burn_rate and policy.name in active:
                active.remove(policy.name)  # recovered; may re-fire later
    return fired


def synthetic_timeline(seed: int = 18, minutes: int = 3 * 24 * 60) -> list[MinuteSample]:
    """Three days of traffic; a scripted bad deploy burns hot for 50 minutes.

    Baseline: ~120 req/min with a healthy 0.05% error rate (below budget).
    Incident: minutes 2000..2049 at 8% errors --- an 80x burn rate, the
    classic "deploy wedged a dependency" shape. Everything seeded, so the
    alert timeline below reproduces on every machine.
    """
    rng = random.Random(seed)
    out: list[MinuteSample] = []
    for t in range(minutes):
        requests = 120 + rng.randint(-10, 10)
        error_rate = 0.08 if 2000 <= t < 2050 else 0.0005
        # errors are a deterministic quota plus seeded jitter around it
        errors = int(requests * error_rate) + (1 if rng.random() < 0.02 else 0)
        out.append(MinuteSample(t, requests, errors))
    return out


def _demo() -> None:
    samples = synthetic_timeline()
    print(f"timeline: {len(samples)} minutes, SLO={SLO}, "
          f"budget={BUDGET_MINUTES_EQUIVALENT:.1f} min-equivalent/30d")
    print(f"{'policy':<10} {'burn':>5} {'long':>6} {'short':>6} {'budget%':>8}")
    for p in POLICIES:
        print(f"{p.name:<10} {p.burn_rate:>5} {p.long_window_min:>5}m "
              f"{p.short_window_min:>5}m {p.budget_consumed_pct:>7.1f}%")
    print("\nalert timeline (minute = minutes since window start):")
    for alert in evaluate(samples):
        hh, mm = divmod(alert.fired_at, 60)
        print(f"  [t={alert.fired_at:>5} ({hh:02d}h{mm:02d}m)] PAGE {alert.policy}"
              f"  long-burn={alert.long_burn:.1f}x short-burn={alert.short_burn:.1f}x")


# ---------------------------------------------------------------------------
# Embedded test suite: pytest burn_rate_sim.py -q
# ---------------------------------------------------------------------------
def test_error_budget_minutes() -> None:
    assert abs(BUDGET_MINUTES_EQUIVALENT - 43.2) < 1e-9  # the famous 43.2 min/month


def test_budget_consumed_percentages_match_workbook_table() -> None:
    assert abs(POLICIES[0].budget_consumed_pct - 2.0) < 1e-9
    assert abs(POLICIES[1].budget_consumed_pct - 5.0) < 1e-9
    assert abs(POLICIES[2].budget_consumed_pct - 10.0) < 1e-9


def test_healthy_traffic_never_pages() -> None:
    samples = synthetic_timeline()
    healthy = [s for s in samples if not 1950 <= s.minute < 2100]
    assert evaluate(healthy) == []


def test_incident_fires_both_pages_but_no_ticket() -> None:
    alerts = evaluate(synthetic_timeline())
    fast = [a for a in alerts if a.policy == "fast-page"]
    assert fast, "an 80x burn for 50 minutes must page the fast policy"
    assert 2005 <= fast[0].fired_at <= 2060  # fires during/right after burst
    # 50 min at 80x still averages 11x over the 6h window -> slow page too
    assert any(a.policy == "slow-page" for a in alerts)
    # but 50 bad minutes out of 3 days stays under the 1x ticket threshold
    assert not any(a.policy == "ticket" for a in alerts)


def test_sustained_low_burn_fires_ticket_only() -> None:
    # 0.3% errors = 3x burn, sustained for the whole timeline
    samples = [
        MinuteSample(t, requests=1000, errors=3) for t in range(3 * 24 * 60 + 400)
    ]
    alerts = evaluate(samples)
    names = {a.policy for a in alerts}
    assert names == {"ticket"}, f"3x burn should page ticket only, got {names}"


def test_window_rate_excludes_boundary_and_handles_empty() -> None:
    samples = [MinuteSample(0, 100, 10), MinuteSample(5, 100, 0)]
    assert window_rate(samples, 5, 5) == 0.0            # (0, 5] holds minute 5 only
    assert window_rate(samples, 5, 6) == 10 / 200       # (-1, 5] holds both minutes
    assert window_rate(samples, 10, 5) == 0.0           # (5, 10] holds neither
    assert window_rate([], 10, 5) == 0.0


if __name__ == "__main__":
    _demo()
\end{minted}

\noindent Its alert timeline for the scripted incident:

\begin{minted}{text}
alert timeline (minute = minutes since window start):
  [t= 2010 (33h30m)] PAGE fast-page  long-burn=14.5x short-burn=76.1x
  [t= 2027 (33h47m)] PAGE slow-page  long-burn=6.2x short-burn=71.7x
\end{minted}

\subsection{Listing 18.6 --- Runbook and Postmortem Generator}
\label{subsec:ch18-listing-templates}

\texttt{incident\_templates.py} turns incident metadata into markdown
skeletons with every mandatory section present: the runbook carries SLOs,
alert entry points, a ten-minute triage checklist, mitigation commands, and
escalation; the postmortem carries the blameless preamble, timeline, impact,
mechanism-level root cause, action items, and an explicit ``lessons for the
SLO'' section. Tests assert the sections exist and that invalid severities
and empty service lists are rejected.

\noindent\textbf{Listing 18.6 --- \texttt{incident\_templates.py}: structure
enforced by code.}
\label{lst:ch18-templates}
\begin{minted}{python}
"""Runbook and blameless-postmortem skeleton generator (stdlib only).

Incident response fails in predictable ways: nobody remembers the rollback
command at 3 a.m., and postmortems devolve into blame or vanish unwritten.
Both are structure problems, and structure is exactly what code is good at.
Given incident metadata, this module emits markdown skeletons with every
mandatory section present and every placeholder explicit, so the human
work is analysis, not formatting.

Run:  python incident_templates.py        (writes demo artifacts)
      pytest incident_templates.py -q     (embedded test suite)
"""

from __future__ import annotations

from dataclasses import dataclass, field
from pathlib import Path

SEVERITIES = {
    "SEV1": "total outage or data loss; all-hands, page immediately",
    "SEV2": "major degradation of a core SLO; page the owning team",
    "SEV3": "minor degradation or single-tenant impact; ticket, next business day",
    "SEV4": "cosmetic or internal-only; backlog",
}


@dataclass(frozen=True)
class ServiceProfile:
    name: str
    slo_target: float
    owner: str
    endpoints: tuple[str, ...]
    dashboards: tuple[str, ...]
    rollback_command: str


@dataclass(frozen=True)
class IncidentMetadata:
    incident_id: str
    title: str
    severity: str
    commander: str
    services: tuple[str, ...]
    started_utc: str
    detected_by: str
    slo_impact: str
    action_items: tuple[str, ...] = field(default=())

    def __post_init__(self) -> None:
        if self.severity not in SEVERITIES:
            raise ValueError(f"severity must be one of {sorted(SEVERITIES)}")
        if not self.services:
            raise ValueError("an incident must name at least one service")


def render_runbook(profile: ServiceProfile) -> str:
    endpoints = "\n".join(f"| `{ep}` | | |" for ep in profile.endpoints)
    boards = "\n".join(f"- {d}" for d in profile.dashboards)
    return f"""# Runbook: {profile.name}

Owner: {profile.owner} | SLO: {profile.slo_target:.3%} availability (30d rolling)

## 1. Service summary
One paragraph: what {profile.name} does, who calls it, what breaks when it
breaks.

## 2. SLOs and alert entry points
| Endpoint | SLI | Alert |
|---|---|---|
{endpoints}

Dashboards:
{boards}

## 3. Triage checklist (first 10 minutes)
1. Confirm the alert is real: check the SLO dashboard, not one log line.
2. Declare severity (see severity table) and open an incident channel.
3. Capture the current trace exemplars before mitigation flushes them.

## 4. Mitigation
- Rollback: `{profile.rollback_command}`
- Feature flag: <flag name and console path>
- Traffic shift / drain: <load balancer procedure>

## 5. Escalation
- Primary on-call -> secondary (15 min ack) -> engineering manager.
"""


def render_postmortem(meta: IncidentMetadata) -> str:
    items = "\n".join(f"- [ ] {a}" for a in meta.action_items) or "- [ ] <action>"
    services = ", ".join(meta.services)
    return f"""# Postmortem: {meta.incident_id} -- {meta.title}

**Severity:** {meta.severity} ({SEVERITIES[meta.severity]})
**Incident commander:** {meta.commander}
**Services:** {services}
**Started (UTC):** {meta.started_utc}
**Detected by:** {meta.detected_by}
**SLO impact:** {meta.slo_impact}

> This document is blameless. We analyze systems, not people: any engineer
> in the same situation, with the same information, might have done the
> same thing. Naming a person as a cause voids the review.

## Summary
Two sentences a VP can read.

## Impact
Who was affected, for how long, in what units (requests failed, SLO budget
burned, revenue or contractual exposure).

## Timeline (all times UTC)
| Time | Event |
|---|---|
| | detection (alert or human?) |
| | mitigation applied |
| | full recovery |

## Root cause
The mechanism, not the culprit: what change, what assumption failed, why
existing guardrails did not catch it.

## What went well / what went poorly
Detection latency, mitigation latency, runbook accuracy, paging quality.

## Action items
{items}

## Lessons for the SLO
Did the alerts fire at the right time? Was the burn-rate table tuned
correctly, or did this incident reveal a gap (add a row, adjust a window)?
"""


def write_artifacts(out_dir: Path, profile: ServiceProfile, meta: IncidentMetadata) -> None:
    out_dir.mkdir(parents=True, exist_ok=True)
    (out_dir / f"runbook_{profile.name}.md").write_text(render_runbook(profile), encoding="utf-8")
    (out_dir / f"postmortem_{meta.incident_id}.md").write_text(
        render_postmortem(meta), encoding="utf-8"
    )


def _demo() -> None:
    profile = ServiceProfile(
        name="fleet-telemetry",
        slo_target=0.999,
        owner="fleet-platform@northwind-robotics.example",
        endpoints=("POST /fleet/telemetry", "GET /fleet/status", "GET /health"),
        dashboards=("SLO / fleet-telemetry (Grafana uid=nr-fleet-slo)",
                    "RED / per-endpoint (Grafana uid=nr-fleet-red)"),
        rollback_command="kubectl rollout undo deploy/fleet-telemetry",
    )
    meta = IncidentMetadata(
        incident_id="NR-2026-0142",
        title="Retry storm after inventory latency regression",
        severity="SEV2",
        commander="a.reyes",
        services=("fleet-gateway", "orders", "inventory"),
        started_utc="2026-08-19T14:03Z",
        detected_by="fast-burn SLO alert (14.4x, 1h/5m)",
        slo_impact="41% of monthly availability budget in 52 minutes",
        action_items=(
            "Cap gateway retries at 2 with jitter (owner: fleet-platform)",
            "Add fan-out panel (requests per client request) to the SLO board",
        ),
    )
    write_artifacts(Path("incident_docs"), profile, meta)
    print(render_postmortem(meta).splitlines()[0])
    print("wrote incident_docs/runbook_fleet-telemetry.md and",
          "incident_docs/postmortem_NR-2026-0142.md")


# ---------------------------------------------------------------------------
# Embedded test suite: pytest incident_templates.py -q
# ---------------------------------------------------------------------------
def _profile() -> ServiceProfile:
    return ServiceProfile(
        name="orders", slo_target=0.9995, owner="orders@northwind-robotics.example",
        endpoints=("POST /orders",), dashboards=("SLO / orders",),
        rollback_command="kubectl rollout undo deploy/orders",
    )


def _meta() -> IncidentMetadata:
    return IncidentMetadata(
        incident_id="NR-2026-0001", title="Test incident", severity="SEV3",
        commander="oncall", services=("orders",), started_utc="2026-01-01T00:00Z",
        detected_by="test", slo_impact="none",
    )


def test_runbook_has_all_mandatory_sections() -> None:
    text = render_runbook(_profile())
    for section in ("Service summary", "SLOs and alert entry points",
                    "Triage checklist", "Mitigation", "Escalation"):
        assert f"## " in text and section in text, f"missing section: {section}"
    assert "kubectl rollout undo deploy/orders" in text
    assert "99.950%" in text  # SLO rendered as a percentage


def test_postmortem_is_blameless_and_complete() -> None:
    text = render_postmortem(_meta())
    for section in ("Summary", "Impact", "Timeline", "Root cause",
                    "Action items", "Lessons for the SLO"):
        assert section in text, f"missing section: {section}"
    assert "blameless" in text.lower()
    assert "SEV3" in text


def test_invalid_severity_rejected() -> None:
    import pytest

    with pytest.raises(ValueError):
        IncidentMetadata(
            incident_id="X", title="t", severity="SEV9", commander="c",
            services=("s",), started_utc="2026-01-01T00:00Z",
            detected_by="test", slo_impact="none",
        )
    with pytest.raises(ValueError):
        IncidentMetadata(
            incident_id="X", title="t", severity="SEV1", commander="c",
            services=(), started_utc="2026-01-01T00:00Z",
            detected_by="test", slo_impact="none",
        )


def test_write_artifacts(tmp_path: Path) -> None:
    write_artifacts(tmp_path, _profile(), _meta())
    assert (tmp_path / "runbook_orders.md").exists()
    assert (tmp_path / "postmortem_NR-2026-0001.md").exists()


if __name__ == "__main__":
    _demo()
\end{minted}
%% ----------------------------------------------------------------------------
\section{Common Pitfalls \& Anti-Patterns}
\label{sec:ch18-pitfalls}

\begin{enumerate}
  \item \textbf{Logging PII, tokens, or request bodies.} One
        \texttt{logger.info("request \%s", body)} and your log backend ---
        with its generous retention, broad ACLs, and vendor UIs --- is now
        a shadow database of everything Chapter~\ref{ch:api_security} told
        you to protect. Log identifiers and outcomes, not payloads; enforce
        it with an allowlist formatter field set and a redaction processor
        in the collector as the second line of defense.
  \item \textbf{High-cardinality metric labels.} \texttt{user\_id},
        \texttt{request\_id}, concrete URL paths, or error-message strings
        as label values are a memory and billing incident on a timer
        (Equation~\eqref{eq:ch18-cardinality}: $90$ series become $4.5$
        million). Rule of thumb: if you would not put the label on a
        dashboard group-by, it does not belong on a metric. Use trace
        attributes and logs for unbounded dimensions.
  \item \textbf{Averaging percentiles, or percentiling averages.} Both are
        statistical category errors (Proposition~\ref{prop:ch18-p99avg}):
        quantiles do not average, and averaging first destroys the tail you
        were measuring. Merge histograms by adding bucket counts; estimate
        percentiles from the merged histogram. If your metrics pipeline
        cannot do that, fix the pipeline before trusting the dashboards.
  \item \textbf{Alerting on every 5xx spike.} A raw error-count page fires
        for a 30-second blip and sleeps through a slow 30-day bleed.
        Alert on budget burn with the window pairs of
        Table~\ref{tab:ch18-burnrate}: fast burn pages, slow burn tickets,
        blips appear on dashboards.
  \item \textbf{Sampling away the errors you need.} A 1\% head sampler
        keeps 1\% of your error traces too --- exactly the traces you will
        beg for during an incident. Sample errors at 100\% (tail sampling
        at the collector, or a keep-all-errors sampler), and aggregate
        metrics \emph{before} sampling, never from sampled spans.
  \item \textbf{Dashboard sprawl without SLO anchoring.} Sixty panels that
        no alert references are a museum. Build dashboards outward from the
        SLO row; any panel that has never been cited in an incident review
        or a runbook is a deletion candidate. Dashboard count is not a
        maturity metric --- decision latency is.
  \item \textbf{Tracing without propagation through queues.} The
        \texttt{traceparent} header dies at the broker unless you
        explicitly inject the context into message metadata on publish and
        extract it on consume (Chapter~\ref{ch:async_apis}). The symptom is
        a forest of single-span traces and a service map with severed
        edges; the fix is one inject/extract pair per messaging boundary.
  \item \textbf{Mean-keyed latency dashboards and SLOs.} ``Average latency
        761\,ms'' sounds acceptable for the distribution of
        Figure~\ref{fig:ch18-longtail} while 1\% of users wait 24 seconds.
        SLOs are threshold ratios (``99\% under 300\,ms''); dashboards show
        p50/p95/p99 from histograms; the mean survives only as a capacity
        planning input.
\end{enumerate}

%% ----------------------------------------------------------------------------
\section{Hands-On Production Lab \& Real-World Case Study}
\label{sec:ch18-lab}

\subsection*{Lab: Propagation, Percentiles, and Pages}

\textbf{Goal.} Reproduce this chapter's three load-bearing claims on your
own machine: trace context survives a three-hop chain, percentiles beat
averages on a long tail, and burn-rate policies page at the mathematically
predicted minutes.

\textbf{Setup.} Save Listings 18.1--18.6 into one directory and create the
virtual environment from Section~\ref{sec:ch18-objectives}.

\begin{enumerate}
  \item \textbf{Prove the wire format.} \texttt{pytest test\_trace\_context.py -q}
        --- 19 passes. Then flip the flags nibble in one valid header to
        \texttt{02} and confirm the parser accepts it but reports
        \texttt{sampled=False}: bit 0 is the only defined bit.
  \item \textbf{Run the instrumented service.} \texttt{python instrumented\_app.py}
        --- three spans export to \texttt{spans.jsonl}, all sharing the
        incoming trace id \texttt{4bf9\ldots4736}, with the parent chain
        asserted. Open \texttt{spans.jsonl} and find the two RED metrics in
        the demo output. Then break propagation: comment out the
        \texttt{PROPAGATOR.inject(outbound)} in \texttt{orders\_create} and
        re-run --- the assertion fails and \texttt{inventory.reserve}
        becomes a root span of a new trace. You have just manufactured the
        most common production tracing bug.
  \item \textbf{Measure the lying average.} \texttt{python latency\_histograms.py}
        --- confirm mean $\approx$ 761\,ms against p99 $\approx$ 23.7\,s,
        then \texttt{pytest latency\_histograms.py -q} (7 passes). Edit
        \texttt{synthetic\_latencies} to drop the tail branch to 1\% and
        watch the mean collapse toward the median while p95 barely moves:
        the mean is a tail detector, not a typical-experience metric.
  \item \textbf{Watch the alerts fire.} \texttt{python burn\_rate\_sim.py}
        --- the fast page fires at $t{=}2010$, the slow page at $t{=}2027$.
        Then edit the incident to a sustained 0.3\% error rate with no
        spike and confirm only the ticket policy fires
        (\texttt{pytest burn\_rate\_sim.py -q} already proves both
        directions: 6 passes). Explain, using
        Equation~\eqref{eq:ch18-burnrate}, why the fast page fired 10
        minutes into the incident and not at its first minute.
  \item \textbf{Generate the paperwork.} \texttt{python incident\_templates.py}
        --- inspect the runbook and postmortem skeletons, then fill the
        postmortem template for the scripted incident as if you had paged
        on it: timeline from the alert log, root cause, and two action
        items.
\end{enumerate}

\subsection{Case Study: The Invisible Retry Storm}
\label{sec:ch18-case}

Northwind Robotics' fleet gateway calls \texttt{orders}, which calls
\texttt{inventory}. On a Thursday deploy, an \texttt{inventory} connection
pool was resized from 50 to 5 (a configuration typo, reviewed and green).
Inventory p50 rose from 22\,ms to 140\,ms and p99 to 1.8\,s --- slow, but
every individual request \emph{succeeded}. The gateway's resilience layer
(Chapter~\ref{ch:microservices}) treated the slowness-plus-occasional-timeout
as transient and retried: two retries with jittered backoff, exactly as
configured.

The logs were a wall of health. Every line was a well-formed request that
completed: 200s everywhere, INFO level, correct correlation ids. Aggregate
error rate: under 0.1\%. Nobody paged on errors, correctly.

What actually happened was visible only in the traces. Spans per trace for
\texttt{POST /fleet/telemetry} jumped from a steady 3 (one per hop) to a
median of 7 --- one gateway span, then \emph{three} \texttt{orders.create}
CLIENT spans where there should be one, each retry fanning into its own
\texttt{inventory.reserve} subtree. The trace fan-out metric --- spans per
trace, per route --- had no dashboard panel, but the trace backend's
service-graph view showed \texttt{orders} receiving $5\times$ its usual
request rate while \emph{client-facing} rate was flat. That mismatch was
the smoking gun: retries multiply downstream load, so the victim service's
traffic diverges from the front door's. The retry storm then deepened the
latency that triggered it --- a self-reinforcing loop that error-rate logs
cannot see because nothing, technically, was failing.

The slow-burn SLO alert (6$\times$, 6h/30m) paged at minute 41 of hour
two: latency SLI burn, not errors. The fix was threefold: the pool typo
reverted; the gateway's retry policy capped at one retry with a per-route
retry \emph{budget} (retries themselves rate-limited,
Chapter~\ref{ch:rate_limiting}); and two permanent additions to the
service dashboard --- a \emph{fan-out panel} (downstream requests per
client request, per route) and a spans-per-trace alert at $2\times$
baseline. The postmortem's headline: \emph{your retry policy is a load
multiplier, and multipliers are invisible to logs.} War-story rules worth
keeping: alert on the topology of traffic (fan-out), not only its volume
and failures; and treat every resilience mechanism as something that can
itself become the incident~\cite{nygard2018releaseit}.

%% ----------------------------------------------------------------------------
\section{Self-Assessment Exercises \& Deep-Dive Questions}
\label{sec:ch18-exercises}

\begin{enumerate}
  \item \textbf{SLI selection.} Northwind Robotics ships a new endpoint,
        \texttt{GET /fleet/route-plan}, that computes asynchronously and
        returns 202 + a polling URL (Chapter~\ref{ch:async_apis}). Write
        SLIs for availability, latency, \emph{and} freshness of the
        finished plan. Which events count as ``good'' for each, and why is
        ``request returned 202'' not an availability SLI for the plan
        itself?
  \item \textbf{Budget arithmetic.} A 99.95\% SLO service burned 14 minutes
        of budget by day 10. What average burn rate does that represent,
        and what is the latest day on which a 2$\times$ sustained burn can
        start without exhausting the month's budget?
  \item \textbf{Window pairs.} Using Equation~\eqref{eq:ch18-burnrate},
        show that a 100\% error burst lasting $t$ minutes raises the 1-hour
        burn rate to $t / (0.001 \cdot 60)$. From that, derive the shortest
        full-outage that fires the 14.4$\times$ fast page, and check your
        answer against Listing~18.5's simulator.
  \item \textbf{Histogram design.} Choose bucket boundaries for an endpoint
        with a 300\,ms latency SLO and explain why at least two buckets
        must straddle 300\,ms. What does the percentile estimate's maximum
        error become at the SLO threshold under your layout?
  \item \textbf{Cardinality audit.} A teammate proposes labels
        \texttt{route}, \texttt{status}, \texttt{tenant\_id} (4{,}000
        tenants), and \texttt{client\_version} (30 versions) on the request
        histogram. Compute the series count, decide which labels survive,
        and say where the rejected dimensions go instead.
  \item \textbf{Propagation through a queue.} Sketch the inject/extract
        code path for telemetry published to a message broker, including
        what happens to the trace when a consumer batches 50 messages into
        one processing span. (Hint: span \emph{links}, not parents.)
  \item \textbf{Fallacy patrol.} Your weekly report shows ``average p99
        latency across instances: 210\,ms.'' List two distinct ways this
        number can understate user pain, and give the correct aggregation.
  \item \textbf{Sampling design.} You may keep 5\% of trace volume.
        Design the sampling policy (head + tail) for a fleet API where
        errors cost you money but slow-and-successful requests cost your
        \emph{customers} money. State precisely which traces are guaranteed
        to survive.
  \item \textbf{Debugging drill.} p99 for \texttt{POST /orders} tripled;
        per-span durations in sampled traces are unchanged, but the gap
        between the SERVER span's start and its first child span grew.
        Name three candidate causes and the USE metric that confirms or
        clears each.
  \item \textbf{Postmortem review.} Given a postmortem whose root-cause
        section reads ``on-call engineer misconfigured the pool,'' rewrite
        it blamelessly: what mechanism questions replace the person
        question, and which action items follow?
\end{enumerate}
%% ----------------------------------------------------------------------------
\section{Interview Q\&A Vault}
\label{sec:ch18-interview}

\subsection*{Junior (Q1--Q10)}
\begin{enumerate}[label=\textbf{Q\arabic*.}, leftmargin=1.6cm]
  \item \textbf{What are the three pillars of observability, and what
        question does each answer?}
        Logs: what happened, once, here (events). Metrics: how often and
        how fast, aggregated cheaply (numbers over time). Traces: where did
        this one request's time go across services (causality). They join
        on shared identifiers --- trace ids in logs, exemplars on metrics.
  \item \textbf{What makes a log line ``structured,'' and why do you care?}
        It is a machine-parseable object (one JSON per line) with a fixed
        schema --- timestamp, level, logger, message, correlation ids ---
        rather than free text. You care because parsing regexes over prose
        during an incident is how minutes of error budget become hours.
  \item \textbf{What may never appear in logs?}
        PII, credentials, tokens, session ids, request/response bodies, and
        anything Chapter~\ref{ch:api_security} classifies as sensitive.
        Logs have wide read access and long retention; treat them as a
        broadcast medium.
  \item \textbf{Define RED and USE.}
        RED --- Rate, Errors, Duration: the service-level view (are users
        hurting?). USE --- Utilization, Saturation, Errors: the
        resource-level view (which pool, queue, or disk is the
        bottleneck?). You page on RED and diagnose with USE.
  \item \textbf{What is a span? A trace?}
        A span is a named, timed operation with attributes, a status, and a
        parent pointer. A trace is all spans sharing one trace id --- the
        tree of one request's work across every service it touched.
  \item \textbf{What does the \texttt{traceparent} header look like?}
        Four dash-separated lowercase-hex fields:
        \texttt{version-traceid-spanid-flags}, e.g.\
        \texttt{00-4bf92f3577b34da6a3ce929d0e0e4736-00f067aa0ba902b7-01}.
        Version \texttt{ff} is invalid; all-zero ids are invalid; flags bit
        0 is \emph{sampled}.
  \item \textbf{What is an SLO, in one sentence?}
        A measurable target on a service level indicator over a rolling
        window --- e.g.\ ``99.9\% of requests succeed over 30 days'' ---
        against which you define an error budget and alerting.
  \item \textbf{How much downtime does a 99.9\% monthly SLO allow?}
        $0.001 \times 43{,}200$ minutes $=$ \textbf{43.2 minutes} per
        month. (99.99\% allows 4.32; 99\% allows 7 h 12 min.)
  \item \textbf{Why is average latency a bad dashboard number?}
        Latency is long-tailed; the mean is dragged into the tail and
        describes no actual user. Use histogram-based p50/p95/p99 ---
        Listing~18.4's demo shows a 761\,ms mean over a 25\,ms median and
        a 23.7\,s p99.
  \item \textbf{What is a correlation id and where does it live?}
        One identifier (here: the W3C trace id) attached to a request at
        the edge and echoed in every log line, span, metric exemplar, and
        problem-details response, so one id retrieves the request's whole
        story.
\end{enumerate}

\subsection*{Mid (Q11--Q20)}
\begin{enumerate}[label=\textbf{Q\arabic*.}, leftmargin=1.6cm, start=11]
  \item \textbf{Walk through the OpenTelemetry architecture.}
        Application code calls the vendor-neutral \emph{API}; the
        \emph{SDK} (configured by the application only) runs samplers, span
        processors, and metric aggregations; \emph{exporters} ship OTLP to
        the \emph{collector}, which batches, redacts, tail-samples, and
        fans out to per-signal backends. Vendor changes stop at collector
        config~\cite{opentelemetrydocs}.
  \item \textbf{Why can't you average p99s across instances or minutes?}
        Quantiles are order statistics, not additive measures; the average
        of per-minute p99s is not any percentile of the underlying
        distribution and systematically understates the tail
        (Proposition~\ref{prop:ch18-p99avg}). Correct path: merge histogram
        bucket counts, then compute the percentile from the merged
        histogram.
  \item \textbf{Explain multi-window multi-burn-rate alerting.}
        Alert on burn rate (observed error rate $\div$ budget rate), and
        fire only when a \emph{long} window (precision: the burn is real
        and sustained) and a \emph{short} window (currency: it is happening
        now) both exceed the threshold --- Equation~\eqref{eq:ch18-mwmb}.
        Canonical 99.9\% table: 14.4$\times$ with 1h+5m (2\% budget), 6$\times$
        with 6h+30m (5\%), 1$\times$ with 3d+6h (10\%, ticket).
  \item \textbf{Head vs tail sampling trade-offs?}
        Head sampling decides at trace start: cheap, unbiased statistics,
        but it discards rare error traces at the sampling rate. Tail
        sampling decides at the collector after buffering whole traces:
        keeps all errors and slow traces, at the cost of collector memory,
        decision latency, and operational complexity. Hybrids head-sample
        the healthy bulk and tail-sample the interesting tail.
  \item \textbf{What is cardinality and why does it bankrupt metrics
        systems?}
        Series count is the product of label cardinalities
        (Equation~\eqref{eq:ch18-cardinality}); each series costs index and
        chunk memory plus ingestion on every scrape. One unbounded label
        (\texttt{user\_id}) turns 90 series into 4.5 million. Labels must
        be dashboard group-bys; unbounded dimensions belong in logs and
        trace attributes.
  \item \textbf{What are exemplars?}
        Trace ids pinned to specific histogram observations. They convert
        ``p99 spiked'' into one clickable slow trace --- the metric-to-
        waterfall bridge --- and they only work because metrics and traces
        share the trace-id namespace.
  \item \textbf{Design SLIs for an asynchronous (202 + polling) API.}
        Three ratios: submission availability (accepted vs.\ rejected
        submissions), completion latency (fraction of jobs finishing within
        $T$ of submission, measured at the poller or job store), and
        freshness/correctness of results where applicable. The HTTP 202
        latency is a client-ergonomics metric, not the job's SLI --- the
        unit of good is the \emph{finished job}, not the acknowledgment.
  \item \textbf{How does \texttt{traceparent} propagate through a queue?}
        It does not --- by itself. There are no HTTP headers on a broker
        message unless you put them there: inject the context into message
        metadata on publish, extract it on consume, and start the consumer
        span as a \texttt{CONSUMER} kind (with span links when batching).
        The OTel propagator API is transport-agnostic; the discipline is
        not automatic.
  \item \textbf{What is the difference between \texttt{tracestate} and
        baggage?}
        \texttt{tracestate}: vendor-specific propagation state (sampling
        decisions, vendor ids), a $\le$32-member \texttt{key=value} list
        with strict syntax, managed by tracers. Baggage: user-defined
        key/value pairs propagated with the trace for application use
        (tenant, plan tier) --- deliberately separate, and a PII leak
        vector if abused because it rides every outbound call.
  \item \textbf{Your p99 latency alert fires. First three moves?}
        (1)~Check the SLO dashboard: is this budget burn or a blip ---
        which endpoint, which burn window? (2)~Pull an exemplar trace from
        the violated histogram bucket and decompose the waterfall:
        dominant span or inter-span gap? (3)~Check USE metrics for the
        implicated resource (pool utilization, queue depth) before reading
        a single log line.
\end{enumerate}

\subsection*{Senior (Q21--Q30)}
\begin{enumerate}[label=\textbf{Q\arabic*.}, leftmargin=1.6cm, start=21]
  \item \textbf{Derive the fast-page window pair from first principles.}
        A 99.9\% SLO allows budget rate $0.001$. To page when 2\% of the
        monthly budget burns in one hour, require burn rate $b$ with
        $b \times 1\text{h} / 720\text{h} = 0.02$, so $b = 14.4$. The 1-hour
        long window estimates the burn with enough samples; the 5-minute
        short window confirms it is still active, so the alert resolves
        after mitigation instead of paging on history.
  \item \textbf{Why do 4xx responses not count against an availability
        SLO?}
        The SLI measures \emph{service} behavior; a 400 means the client
        sent an invalid request --- counting it would let one buggy
        integrator burn your budget and page you for their bug. Track 4xx
        as a separate indicator (sudden 4xx shifts do reveal bad deploys
        and doc drift), but the availability SLI counts 5xx and transport
        failures.
  \item \textbf{Your traces show healthy spans but the gaps between spans
        grew. Where did the time go?}
        Between-span time is work nobody instrumented: connection
        acquisition from an exhausted pool, TLS handshakes on cold
        connections, DNS, thread/executor queueing, GC pauses, or
        client-side retries (check fan-out first --- see
        Section~\ref{sec:ch18-case}). Confirm with USE metrics on the
        suspected resource, then add spans for the confirmed stage.
  \item \textbf{How do you keep metrics correct while sampling traces?}
        Aggregate metrics from every request \emph{before} sampling ---
        counters and histograms are unconditional; only span export is
        sampled. Never derive request rates or latency percentiles from
        sampled spans: your RED metrics would be off by the sampling
        factor, silently.
  \item \textbf{Design a latency SLO for an endpoint with a heavy tail of
        legitimate work (large fleet exports).}
        Split the SLI by request class (separate export vs.\ interactive
        routes, different thresholds), pick the threshold from the
        histogram at the knee rather than from a wish, and express the SLO
        as a threshold ratio (``99\% of interactive requests under
        300\,ms''). If one route genuinely has two populations, one SLI
        will lie to one of them.
  \item \textbf{What breaks when you attach \texttt{http.route} as the raw
        path instead of the template?}
        \texttt{/fleet/robot-4821} instead of \texttt{/fleet/\{id\}} makes
        every id a new label value: per-request series creation, index
        blowup, dashboards that cannot aggregate, and eventually a metrics
        backend that drops or bills you into noticing. Route templates are
        a cardinality control, not a cosmetic choice.
  \item \textbf{Explain the alert-fatigue failure mode and its fix.}
        Pages that do not require immediate action train the on-call to
        ignore pages; response latency to \emph{real} pages degrades ---
        fatigue is a system output, not a character flaw. Fix: page only on
        symptom-level budget burn (fast/slow policies), ticket the rest,
        delete or demote alerts that fire unactionably, and review page
        volume as a budget.
  \item \textbf{How would you verify a third-party client's claim that
        ``your API was down for us at 14:03''?}
        Ask for the correlation id or \texttt{traceparent} they sent (you
        taught them to, via problem-details \texttt{traceId}); query traces
        and logs for that id; then check the SLI window around 14:03 for
        their route and region. If their requests never appear in server
        telemetry, the failure is between them and your edge --- which is
        also a finding.
  \item \textbf{When is a gauge the wrong instrument?}
        For anything you will want to aggregate across instances or
        percentile over time: gauges capture instantaneous values at scrape
        time and lose everything between scrapes. Rates come from counters,
        distributions from histograms; gauges are for levels that are
        meaningful instantaneously (queue depth, pool in-use).
  \item \textbf{What belongs in a runbook entry for one alert?}
        The SLO it protects; the dashboard that confirms it; the triage
        checklist (ten minutes, ordered); mitigation with exact commands
        (rollback, flag, drain); escalation path with ack timeouts; and the
        date it was last exercised. A runbook that has never been drilled
        is a hypothesis.
\end{enumerate}

\subsection*{Staff (Q31--Q36)}
\begin{enumerate}[label=\textbf{Q\arabic*.}, leftmargin=1.6cm, start=31]
  \item \textbf{Design the observability architecture for a 40-service API
        platform. What do you standardize centrally vs.\ leave to teams?}
        Central: the collector tier (endpoints, redaction, tail sampling,
        per-tenant quotas), the trace-context and baggage policy, metric
        naming and label vocabularies (route templates, status classes),
        the SLO framework and burn-rate alerting library, exemplar wiring,
        and the dashboard/slo-as-code templates. Teams: SLI thresholds and
        SLO targets for their services, domain attributes on spans, and
        their runbooks. The invariant: propagation and cardinality
        discipline are platform concerns; target-setting is an ownership
        concern.
  \item \textbf{When would you deliberately \emph{lower} an SLO?}
        When the budget is consistently unspent and users cannot perceive
        the marginal reliability --- an over-tight SLO is a tax on velocity
        (frozen releases, rejected experiments) buying nothing. SLOs are
        set from user pain and business tolerance, not from what the
        architecture happens to achieve; ratcheting to the observed maximum
        is the stealth failure mode.
  \item \textbf{How do error budgets change organizational behavior, and
        how do you keep them honest?}
        The budget converts reliability from a vibe into a shared currency:
        unspent budget authorizes risk (faster releases, experiments);
        exhausted budget shifts the roadmap to reliability by policy, not
        by argument~\cite{beyer2016sre}. Keeping it honest: the SLI must be
        user-visible (no internal-health proxies), the measurement window
        and query are versioned and reviewed, and budget policy changes are
        signed off like contracts --- otherwise teams relitigate the ruler
        instead of fixing the service.
  \item \textbf{Critique: ``we sample 10\% of traces and compute our SLO
        from the sampled spans.''}
        Two independent errors. First, metrics from sampled spans are
        biased by the sampler (and catastrophically so under tail sampling,
        which \emph{over-represents} errors and slow traces --- your SLI
        would read worse than reality, or better under head sampling if
        errors correlate with sampling bugs). Second, SLOs must be computed
        from complete event data: unsampled request logs, load-balancer
        records, or pre-sampling metric aggregations. Traces are for
        diagnosis; SLIs are for accounting; never mix the supplies.
  \item \textbf{A queue-based pipeline drops trace continuity at three
        broker boundaries. Design the fix and its failure modes.}
        Inject \texttt{traceparent}/\texttt{tracestate} into message
        metadata at every publish (one carrier-agnostic propagator call),
        extract at every consume, start \texttt{PRODUCER}/\texttt{CONSUMER}
        spans, and use span \emph{links} for fan-in batches so one
        consumer span can reference fifty parents without lying about
        causality. Failure modes to design for: metadata size limits on the
        broker, consumers on old library versions that drop unknown fields
        (contract-test the metadata schema), and the temptation to log the
        baggage --- which now travels everywhere.
  \item \textbf{Your CEO asks: ``why did we invest in all this telemetry
        when the outage still lasted 40 minutes?'' Defend the program with
        numbers.}
        Compare against the counterfactual the telemetry itself records:
        detection time (alert fired at minute 4 of burn, not at a customer
        ticket at minute 40), diagnosis time (exemplar-to-root-cause in
        minutes via the waterfall vs.\ hours of log archaeology), and
        budget math (40 minutes consumed vs.\ 43.2 allowed --- the
        month survived). Observability does not prevent incidents; it
        compresses detection and diagnosis, which are the two phases
        engineering actually controls. Then say what will shrink
        mitigation: better runbooks and rehearsed rollbacks --- the program
        is how you know that.
\end{enumerate}
%% ----------------------------------------------------------------------------
\section{External Bibliography \& Further Reading}
\label{sec:ch18-bibliography}

\begin{itemize}
  \item B.~Beyer, C.~Jones, J.~Petoff, N.~Murphy (eds.), \emph{Site
        Reliability Engineering} (O'Reilly, 2016) --- the source of SLOs,
        error budgets, and the elimination-of-toil philosophy this chapter
        operationalizes~\cite{beyer2016sre}.
  \item B.~Beyer, N.~Murphy, D.~Rensin, K.~Kawahara, S.~Thorne (eds.),
        \emph{The Site Reliability Workbook} (O'Reilly, 2018) --- chapter~4
        (``Implementing SLOs'') and chapter~5 (``Alerting on SLOs'') are
        the origin of the multi-window multi-burn-rate table implemented in
        Listing~18.5 (uncited in text; the practical companion to the SRE
        book).
  \item OpenTelemetry project documentation --- API/SDK concepts, the
        collector, and the HTTP semantic conventions referenced throughout
        Section~\ref{sec:ch18-mechanics}~\cite{opentelemetrydocs}.
  \item W3C, \emph{Trace Context} (Recommendation) and \emph{Baggage} ---
        the normative specifications for \texttt{traceparent} and
        \texttt{tracestate} that Listings 18.1--18.2 implement by hand
        (uncited in text; short, readable, and worth reading in full).
  \item C.~Majors, L.~Fong-Jones, G.~Miranda, \emph{Observability
        Engineering} (O'Reilly, 2022) --- the high-cardinality,
        high-dimensionality argument for event-first observability and the
        clearest treatment of exemplar-style debugging loops (uncited in
        text).
  \item Prometheus documentation, \emph{Histograms and Summaries} --- why
        percentiles must come from histogram buckets, the
        \texttt{histogram\_quantile} interpolation behind
        Equation~\eqref{eq:ch18-histpct}, and the unaggregability of
        summary quantiles (uncited in text).
  \item J.~Turnbull, \emph{The Art of Monitoring} --- the opinionated
        end-to-end framing of metrics, alerting, and dashboard design that
        shaped Section~\ref{subsec:ch18-ops} (uncited in text).
  \item RFC~9110, \emph{HTTP Semantics} --- status-code classes that define
        ``good'' versus ``bad'' events in availability SLIs, and the header
        semantics trace context rides on~\cite{rfc9110}.
  \item M.~Kleppmann, \emph{Designing Data-Intensive Applications} --- the
        messaging and stream-processing background for the
        propagation-through-queues discussion~\cite{kleppmann2017ddia}.
  \item M.~Nygard, \emph{Release It!}, 2nd ed. --- the stability patterns
        whose interactions with telemetry produced
        Section~\ref{sec:ch18-case}'s retry storm~\cite{nygard2018releaseit}.
  \item FastAPI documentation --- middleware and the ASGI surface
        instrumented in Listing~18.3~\cite{fastapidocs}.
  \item A.~H\'{e}tu, \emph{The HdrHistogram} project documentation --- the
        log-bucket layout and its bounded-relative-error guarantee
        re-implemented in Listing~18.4 (uncited in text).
\end{itemize}

%% ----------------------------------------------------------------------------
\section{Capstone Projects}
\label{sec:ch18-capstones}

\subsection*{Capstone 18.1 --- Trace Context Conformance Pack \hfill [junior]}
\textbf{Objective:} Prove you own the wire format by extending the parser
into a conformance checker.

\textbf{Inputs/Datasets:} Listings 18.1--18.2; a hand-written file of 20
\texttt{traceparent} headers (10 valid, 10 invalid in distinct ways) plus
5 \texttt{tracestate} headers.

\textbf{Steps:} (1)~Write the header corpus as a data file, one header per
line with an expected verdict. (2)~Build a checker that runs every corpus
line through \texttt{parse\_traceparent} and reports pass/fail per case.
(3)~Add version-\texttt{01} headers with trailing fields and verify
forward-compatible parsing. (4)~Add a property test: for 500 seeded random
contexts, \texttt{parse(header(x)) == x}.

\textbf{Acceptance Criteria:}
\begin{itemize}
  \item[$\square$] All 25 corpus cases produce the expected verdict.
  \item[$\square$] The seeded round-trip property test passes 500/500.
  \item[$\square$] Every rejection message names the offending field.
\end{itemize}

\textbf{Stretch Goals:} Fuzz with random truncations and case flips; assert
the parser never throws anything but \texttt{ValueError}.

\subsection*{Capstone 18.2 --- RED Metrics for a Real Route Set \hfill [mid]}
\textbf{Objective:} Extend Listing~18.3 into a per-endpoint RED story for
Northwind Robotics' four public routes.

\textbf{Inputs/Datasets:} Listing~18.3; routes \texttt{POST /fleet/telemetry},
\texttt{GET /fleet/status}, \texttt{GET /fleet/\{id\}}, \texttt{GET /health}.

\textbf{Steps:} (1)~Record counter and histogram with attributes
\texttt{route} (template only), \texttt{status\_class}
(\texttt{2xx}/\texttt{4xx}/\texttt{5xx}), nothing else. (2)~Drive 2{,}000
seeded synthetic requests with a per-route latency model and a 0.2\% error
rate on one route. (3)~Read back the \texttt{InMemoryMetricReader} and
render a per-route RED table (rate, error ratio, p50/p99 from bucket
counts). (4)~Prove the cardinality bound: assert the exported series count
equals routes $\times$ status classes, exactly.

\textbf{Acceptance Criteria:}
\begin{itemize}
  \item[$\square$] RED table prints with per-route error ratios matching
        the seeded model within sampling jitter.
  \item[$\square$] Series-count assertion passes; adding a
        \texttt{user\_id} label makes it fail loudly (demonstrate, then
        revert).
  \item[$\square$] Everything deterministic under one seed; no network.
\end{itemize}

\textbf{Stretch Goals:} Attach exemplars (one trace id per histogram bucket
observation) and print one exemplar id for the p99 bucket.

\subsection*{Capstone 18.3 --- SLO Workbench \hfill [mid]}
\textbf{Objective:} Turn Listing~18.5 into a policy design tool and defend
one design change with data.

\textbf{Inputs/Datasets:} Listing~18.5; three scripted incidents you write
(a 3-minute total outage, a 6-hour 2\% error bleed, a flapping
30-seconds-on/30-seconds-off failure).

\textbf{Steps:} (1)~Encode the three incidents as timeline generators with
distinct seeds. (2)~Run the stock policy table against each; record
time-to-page and false positives. (3)~Propose one modification (e.g.\ a
third window pair at 3$\times$/1d+2h, or a flapping-damping reset rule) and
show its effect on all three incidents. (4)~Write a one-page memo: which
table ships, and why, in budget-consumption terms.

\textbf{Acceptance Criteria:}
\begin{itemize}
  \item[$\square$] All three incidents produce a documented, reproducible
        alert timeline.
  \item[$\square$] The flapping incident is shown to reset-and-refire under
        the stock table, and your modification addresses it.
  \item[$\square$] The memo states budget math for every policy row.
\end{itemize}

\textbf{Stretch Goals:} Add alert-resolution tracking (when would each page
auto-resolve?) and report mean time to resolve per policy.

\subsection*{Capstone 18.4 --- Queue Propagation End to End \hfill [senior]}
\textbf{Objective:} Close the tracing gap this chapter warns about: one
trace across HTTP \emph{and} a message broker.

\textbf{Inputs/Datasets:} Listing~18.3; an in-process fake broker you write
(dict of queues; no network, deterministic).

\textbf{Steps:} (1)~Publish telemetry from the gateway to the fake broker,
injecting \texttt{traceparent}/\texttt{tracestate} into message metadata.
(2)~A consumer extracts context and starts a \texttt{CONSUMER} span parented
correctly; a batch consumer uses span links for its 50 parents. (3)~Export
and assert: one trace id spans HTTP ingress, publish, and consume; batch
consumption produces one span with 50 links. (4)~Write the failure-mode
test: a legacy consumer that drops unknown metadata fields produces a
detectably orphaned span (assert and log a warning).

\textbf{Acceptance Criteria:}
\begin{itemize}
  \item[$\square$] The exported tree crosses the broker boundary with one
        trace id and correct parentage.
  \item[$\square$] Batch consumption uses links, not a false single parent.
  \item[$\square$] The legacy-consumer test detects the orphaned span
        programmatically.
\end{itemize}

\textbf{Stretch Goals:} Add a redaction pass that strips baggage at the
queue boundary and prove no baggage key survives into consumer spans.

\subsection*{Capstone 18.5 --- Incident Program in a Box \hfill [staff]}
\textbf{Objective:} Stand up the complete operational loop for the
fleet-telemetry platform and run a game day against it.

\textbf{Inputs/Datasets:} Listings 18.3--18.6; the case study of
Section~\ref{sec:ch18-case} as the game-day scenario.

\textbf{Steps:} (1)~Define SLOs (availability + latency + freshness) for
three services with per-endpoint thresholds. (2)~Generate runbooks from
Listing~18.6 for all three services and fill every placeholder.
(3)~Script the retry-storm incident into the simulator's timeline plus a
fan-out metric; verify the fast/slow pages and show that a
spans-per-trace alert fires \emph{before} the latency burn pages.
(4)~Run the game day: one teammate injects the fault, another follows only
the runbook; measure detection, mitigation, and resolution times.
(5)~Produce the blameless postmortem from the generator, including the
``lessons for the SLO'' section with at least one concrete alerting change.

\textbf{Acceptance Criteria:}
\begin{itemize}
  \item[$\square$] SLOs, runbooks, and dashboards-as-tables exist for all
        three services and cite one another.
  \item[$\square$] The game day produces measured detection/mitigation/
        resolution times, all reproducible from seeds.
  \item[$\square$] The postmortem is complete, blameless, and its action
        items include one alerting-topology change with an owner.
  \item[$\square$] A second run by a different pair reproduces the
        timelines within one minute of simulated time.
\end{itemize}

\textbf{Stretch Goals:} Add a budget-policy gate: when the simulator shows
$>$75\% budget consumed, the deploy script refuses to ship and explains why.

%% ----------------------------------------------------------------------------
\section{Python Dependencies Appendix}
\label{sec:ch18-deps}

All code runs offline after a one-time install. From PowerShell:

\begin{minted}{text}
python -m venv .venv
.\.venv\Scripts\Activate.ps1
pip install -r requirements.txt
\end{minted}

\noindent\textbf{requirements.txt:}
\begin{minted}{text}
opentelemetry-sdk>=1.25
opentelemetry-instrumentation-fastapi>=0.46b0
fastapi>=0.110
pytest>=8
\end{minted}

For real deployments (never needed by this chapter's listings): add the
OTLP exporter package (\texttt{opentelemetry-exporter-otlp}) and point it
at a collector endpoint --- the \texttt{SpanExporter} interface Listing~18.3
implements is the same one OTLP satisfies, so the swap is constructor-deep.

\subsection{opentelemetry-sdk ($\geq$ 1.25)}
\textbf{Description:} The OpenTelemetry SDK for Python: tracer and meter
providers, span processors, samplers, metric views, and the propagator API.
\textbf{Why included:} it is the entire telemetry backbone of Listing~18.3
--- and its \texttt{SpanExporter} interface is what lets the chapter replace
OTLP with an offline JSON-lines exporter without touching application
code. \textbf{Alternatives:} the Prometheus \texttt{prometheus\_client}
(metrics only --- no traces, no logs, and a pull model that changes your
deployment shape); the Datadog SDK / \texttt{ddtrace} (vendor-locked but
zero-collector); \texttt{structlog} for structured logging specifically
(excellent, and composable with the stdlib handler used here --- but it is
a logging library, not a telemetry pipeline). \textbf{Installation:}
\texttt{pip install "opentelemetry-sdk>=1.25"}.

\subsection{opentelemetry-instrumentation-fastapi ($\geq$ 0.46b0)}
\textbf{Description:} Auto-instrumentation that wraps every FastAPI route
in a SERVER span with semantic-convention attributes and extracts incoming
\texttt{traceparent} headers. \textbf{Why included:} it is the difference
between teaching propagation as a diagram and having it happen on every
request for free; Listing~18.3 instruments the app with it and then
demonstrates the manual inject/extract path it automates.
\textbf{Alternatives:} hand-written ASGI middleware (full control, full
maintenance); the \texttt{opentelemetry-instrumentation-asgi} package it
builds on (framework-agnostic, more manual); vendor agents (locked).
\textbf{Installation:}
\texttt{pip install "opentelemetry-instrumentation-fastapi>=0.46b0"} ---
the pre-1.0 instrumentation packages version in lockstep with the SDK
(0.46b0 pairs with SDK 1.25).

\subsection{fastapi ($\geq$ 0.110)}
\textbf{Description:} The ASGI framework hosting the instrumented service.
\textbf{Why included:} the chapter's service is a FastAPI app because the
book's services are FastAPI apps (Chapter~\ref{ch:fastapi}); the framework
supplies routing, middleware, and the request objects the instrumentation
wraps~\cite{fastapidocs}. \textbf{Alternatives:} Starlette directly (the
instrumentation targets ASGI and would work unchanged); Flask via
\texttt{opentelemetry-instrumentation-flask} (WSGI; identical concepts);
a gateway-level service mesh (sidecar telemetry without code changes, at
the price of losing in-process spans). \textbf{Installation:}
\texttt{pip install "fastapi>=0.110"}.

\subsection{pytest ($\geq$ 8)}
\textbf{Description:} Test framework running every embedded suite in this
chapter (36 tests across four files). \textbf{Why included:} the chapter's
claims are executable --- the header round-trips, histogram percentile
proofs, and burn-rate policy tests all run under it.
\textbf{Alternatives:} \texttt{unittest} (stdlib; the suites would survive
a port, but \texttt{parametrize} and \texttt{tmp\_path} would need
hand-rolling); \texttt{nose2} (unmaintained momentum). \textbf{Installation:}
\texttt{pip install "pytest>=8"}.

\subsection*{Ecosystem alternatives worth knowing}
For production telemetry beyond this chapter's offline exporters: the
OpenTelemetry Collector itself (the binary this chapter simulates),
Prometheus + \texttt{prometheus\_client} for metrics-only deployments,
Grafana Tempo/Jaeger for trace storage, Loki or Elasticsearch for logs,
and \texttt{structlog} when you want richer structured logging pipelines
than a stdlib formatter. Evaluate each against
Section~\ref{sec:ch18-mechanics}: which signal, which cardinality story,
which sampling point, which failure mode when the backend is down.

%% ----------------------------------------------------------------------------
\section{Chapter Summary \& Key Takeaways}
\label{sec:ch18-summary}

\begin{itemize}
  \item The three pillars answer different questions and join on one
        identifier: structured JSON logs (levels, sampling, never PII or
        secrets), RED metrics for services and USE for resources, and
        traces as parent-linked span trees propagated by context.
  \item OpenTelemetry's layering is the architecture: code to the API,
        configure the SDK, export to a collector, and let the collector own
        vendor coupling --- the offline file exporter in Listing~18.3 is
        the same interface OTLP satisfies~\cite{opentelemetrydocs}.
  \item \texttt{traceparent} is four hex fields --- version, trace id,
        span id, flags --- with all-zero ids and version \texttt{ff}
        invalid and bit 0 meaning \emph{sampled}
        (Table~\ref{tab:ch18-traceparent}); queues do not propagate it for
        you --- inject and extract at every messaging boundary.
  \item Latency is a histogram. Percentiles come from bucket counts
        (Equation~\eqref{eq:ch18-histpct}); averaging first destroys the
        tail (Proposition~\ref{prop:ch18-p99avg}); cardinality is
        multiplicative and priced per series
        (Equation~\eqref{eq:ch18-cardinality}); exemplars bridge metrics to
        traces.
  \item SLOs turn telemetry into judgment: availability, latency, and
        freshness SLIs as good/total ratios; error budgets (43.2 minutes
        per month at 99.9\%); and multi-window multi-burn-rate alerts ---
        14.4$\times$ over 1h+5m, 6$\times$ over 6h+30m, 1$\times$ over
        3d+6h --- that page on budget burn instead of noise
        (Table~\ref{tab:ch18-burnrate})~\cite{beyer2016sre}.
  \item Operations is a discipline, not a mood: correlation ids end to end,
        dashboards anchored to SLOs, latency debugged by timeline
        decomposition (span or gap first, logs last), incidents run with
        severities, roles, runbooks, and blameless postmortems.
  \item The retry storm of Section~\ref{sec:ch18-case} is the exam
        question: resilience mechanisms multiply load invisibly to logs;
        alert on the topology of traffic (fan-out, spans per trace), not
        only its volume and errors.
\end{itemize}

%% ----------------------------------------------------------------------------
\section{Quick-Reference Card}
\label{sec:ch18-quickref}

\subsection*{SLI Formulas}
Availability $= \frac{\#\,\text{status} < 500}{\#\,\text{requests}}$
$\bullet$ Latency $= \frac{\#\,\text{requests} < T}{\#\,\text{requests}}$
(from histogram buckets) $\bullet$ Freshness $= \frac{\#\,\text{reads newer
than } D}{\#\,\text{reads}}$ $\bullet$ Error budget $= (1-\text{SLO}) \times
\text{window}$: 99.9\%/30d $=$ 43.2 min; 99.95\% $=$ 21.6 min; 99.99\% $=$
4.32 min $\bullet$ Burn rate $= \frac{\text{error rate}}{1 - \text{SLO}}$
$\bullet$ Budget consumed $=$ burn $\times\, w / 43{,}200$ min.

\subsection*{Burn-Rate Window Table (99.9\% SLO)}
Fast page: 14.4$\times$, long 1h + short 5m, 2\% budget $\bullet$ slow
page: 6$\times$, long 6h + short 30m, 5\% budget $\bullet$ ticket:
1$\times$, long 3d + short 6h, 10\% budget $\bullet$ fire $\iff$
burn$_{\text{long}} \ge b \land$ burn$_{\text{short}} \ge b$; reset when
the long window clears.

\subsection*{\texttt{traceparent} Field Reference}
\texttt{version-traceid-spanid-flags} $\bullet$ version: 2 hex,
\texttt{ff} forbidden, \texttt{00} forbids extra fields $\bullet$ trace id:
32 hex, nonzero, constant across hops $\bullet$ span id: 16 hex, nonzero,
fresh per hop, recorded as the callee's parent $\bullet$ flags: 2 hex, bit
\texttt{0x01} $=$ sampled; echo only defined bits $\bullet$
\texttt{tracestate}: $\le$32 vendor \texttt{key=value} members $\bullet$
baggage: user key/values, propagates everywhere --- never PII.

\subsection*{Field Discipline (compressed)}
Logs: JSON, levels mean something, sample hot paths, no PII/tokens/bodies
$\bullet$ Metrics: RED per route template + status class only; histograms
for latency; exemplars on; series $=$ product of label cardinalities
$\bullet$ Traces: SERVER/CLIENT/INTERNAL/PRODUCER/CONSUMER kinds; metrics
aggregated \emph{before} sampling; errors kept at 100\% $\bullet$ Debug
order: SLO dashboard $\to$ exemplar waterfall $\to$ dominant span or gap
$\to$ USE metrics $\to$ correlated logs $\bullet$ Incidents: severity with
a response contract, commander coordinates (never debugs), runbook first
ten minutes, postmortem blames mechanisms.

\section{Standardized Evidence Addendum}
\subsection{Benchmark Snapshot Template}
\begin{table}[h]
  \centering
  \small
  \begin{tabularx}{\textwidth}{l c c c c}
    \toprule
    \textbf{Config} & \textbf{p50 latency} & \textbf{p99 latency} & \textbf{Error rate} & \textbf{Throughput} \\
    \midrule
    Baseline & -- & -- & -- & 1.00x \\
    Candidate A & -- & -- & -- & -- \\
    Candidate B & -- & -- & -- & -- \\
    \bottomrule
  \end{tabularx}
  \caption{Minimum benchmark table to keep per chapter.}
\end{table}

\subsection{Request Trace Template}
\begin{itemize}
  \item \textbf{Request:} one representative API call for this chapter's topic.
  \item \textbf{Before:} behavior (headers, payload, latency, failure mode) with the naive approach.
  \item \textbf{After:} behavior with the technique in this chapter.
  \item \textbf{Result:} what changed in correctness, resilience, latency, and operability.
\end{itemize}

\subsection{Cost and Latency Budget Snapshot}
\begin{table}[h]
  \centering
  \small
  \begin{tabularx}{\textwidth}{l c c}
    \toprule
    \textbf{Stage} & \textbf{Latency budget} & \textbf{Cost driver} \\
    \midrule
    TLS / connection setup & -- & handshake round trips, keep-alive hit rate \\
    Request validation & -- & schema complexity, CPU per request \\
    Business logic and I/O & -- & downstream fan-out, query cost \\
    Serialization and egress & -- & payload size, compression ratio \\
    \bottomrule
  \end{tabularx}
  \caption{Operational budget template for this chapter's technique.}
\end{table}

\section{Configuration Preset Template}
\begin{minted}{text}
# api_policy.yaml
policy_version: "v1.0.0"
component: "<chapter_topic>"
timeout_ms: 0
retry_budget: 0
fallback_strategy: "fail_fast"
rollback_trigger: "error_budget_burn_gt_threshold"
\end{minted}

\section{CI Quality Gate Specification}
\begin{itemize}
  \item contract tests green against the pinned OpenAPI/AsyncAPI/Protobuf spec,
  \item no breaking schema changes without a version bump,
  \item error responses conform to RFC 9457 problem-details shape,
  \item benchmark deltas within approved regression bounds (latency and throughput),
  \item policy version and rollback plan present in release metadata.
\end{itemize}

\section{Incident Runbook Template}
\begin{enumerate}
  \item Detect regression from SLO burn-rate alerts or consumer reports.
  \item Scope impacted endpoints, versions, and consumer segments.
  \item Contain impact (rollback, feature flag, rate-limit tightening, or traffic shift).
  \item Diagnose with traces, structured logs, and contract diffs.
  \item Recover with validated fix and verified rollout procedure.
  \item Prevent recurrence via new regression tests and CI gates.
\end{enumerate}

\section{Cross-Chapter Decision Card}
\begin{table}[h]
  \centering
  \small
  \begin{tabularx}{\textwidth}{l X X}
    \toprule
    \textbf{Dimension} & \textbf{Choose this approach when} & \textbf{Prefer an alternative when} \\
    \midrule
    Use case & -- & -- \\
    Latency budget & -- & -- \\
    Operational cost & -- & -- \\
    Team maturity & -- & -- \\
    \bottomrule
  \end{tabularx}
  \caption{Decision card to complete per chapter.}
\end{table}

\section{ROI Model and Build-vs-Buy}
\begin{itemize}
  \item \textbf{Build when:} the capability is differentiating, compliance forbids vendors, or scale makes vendor pricing irrational.
  \item \textbf{Buy when:} the capability is commodity (auth, gateways, observability), time-to-market dominates, or staffing is thin.
  \item \textbf{Hidden costs to model:} on-call load, upgrade cadence, expertise ramp, data egress fees.
\end{itemize}

\section{Disaster Recovery Notes (RPO/RTO)}
\begin{itemize}
  \item Define recovery point objective (RPO): maximum acceptable data loss window.
  \item Define recovery time objective (RTO): maximum acceptable restoration time.
  \item Test restores regularly; an untested backup is not a backup.
\end{itemize}

